% 子文件：A_统计图_小学
% 本文件包含 17 个 section



\newpage

\section{解答：数据的初步收集与统计图绘制}

\vspace{0.4cm}

\noindent\textbf{3. 根据第2题中收集的数据，分别完成下面的统计图。}

\vspace{0.6cm}

\begin{center}
\textbf{\large 小华一至六年级体重变化情况统计图}

\vspace{0.4cm}
\makebox[9.8cm][r]{2026年1月}

\vspace{0cm}
% =====================================================================
% 【折线图绘制思路】
% 1. X 轴分配 7 列网格空间，第一条纵线映射为“一年级”位置 (x=1)。
% 2. Y 轴映射机制：由于大网格是 5kg 一级，我们将全系统 Y 单位等比例压缩在 0.5cm，通过除以 5 (例如 24/5 = 4.8) 进行精确定位。
% 3. 各点使用标准的 \filldraw 描黑并联通成线，数据签注定位于节点右下方以免干涉连接线势。
% =====================================================================
\begin{tikzpicture}[x=1.4cm, y=0.5cm, font=\small]
  % 绘制网格
  \draw[step=1, black, thin] (0,0) grid (7,8);
  
  % Y 轴标签
  \node[above left] at (0, 8.2) {体重/kg};
  \foreach \y/\val in {0/0, 1/5, 2/10, 3/15, 4/20, 5/25, 6/30, 7/35, 8/40} {
    \node[left] at (0, \y) {\val};
  }
  
  % X 轴标签
  \node[below] at (1, 0) {一年级};
  \node[below] at (2, 0) {二年级};
  \node[below] at (3, 0) {三年级};
  \node[below] at (4, 0) {四年级};
  \node[below] at (5, 0) {五年级};
  \node[below] at (6, 0) {六年级};
  \node[below] at (7, 0) {年级};

  % 折线与数据点
  \draw[thick, black] (1, 4.8) -- (2, 5.2) -- (3, 5.6) -- (4, 6.2) -- (5, 6.8) -- (6, 7.6);
  
  \filldraw[black] (1, 4.8) circle (1.8pt) node[below right, inner sep=2pt] {24};
  \filldraw[black] (2, 5.2) circle (1.8pt) node[below right, inner sep=2pt] {26};
  \filldraw[black] (3, 5.6) circle (1.8pt) node[below right, inner sep=2pt] {28};
  \filldraw[black] (4, 6.2) circle (1.8pt) node[below right, inner sep=2pt] {31};
  \filldraw[black] (5, 6.8) circle (1.8pt) node[below right, inner sep=2pt] {34};
  \filldraw[black] (6, 7.6) circle (1.8pt) node[below right, inner sep=2pt] {38};

\end{tikzpicture}
\end{center}

\vspace{0.8cm}

\begin{center}
\textbf{\large 六年级一班同学家庭成员人数情况统计图}

\vspace{0.4cm}
\makebox[9cm][r]{2026年1月}

\vspace{0cm}
% =====================================================================
% 【条形图绘制思路】
% 1. X 轴原样扩展为 9 列网格区块（匹配由原始 10 条竖线切割出的 9x5 方格面）。
% 2. “空一格画一格”排版：分别在 [1,2]、[3,4]、[5,6]、[7,8] 跨度区间内部，满格绘制条形柱体，并给下方中间位置标注归属类群。
% 3. 映射关系：1 大格代理 5 人，转换参数纵向通过 值/5 来获取 (例如 18人对应的画高为 3.6 单元)。
% 4. 条柱配置：采用满格紧贴两侧竖线 (Rectangle x到x+1) 的绘制法，教辅常配的浅灰色 (gray!30) 内填搭配黑色边框，数字附挂于项端。
% =====================================================================
\begin{tikzpicture}[x=1cm, y=0.8cm, font=\small]
  % 网格
  \draw[step=1, black, thin] (0,0) grid (9,5);
  
  % Y 轴标签
  \node[above left] at (0, 5.2) {家庭数};
  \foreach \y/\val in {0/0, 1/5, 2/10, 3/15, 4/20, 5/25} {
    \node[left] at (0, \y) {\val};
  }
  
  % X 轴标签
  \node[below] at (1.5, 0) {2人};
  \node[below] at (3.5, 0) {3人};
  \node[below] at (5.5, 0) {4人};
  \node[below] at (7.5, 0) {5人};
  \node[below] at (9, 0) {成员人数};

  % 条形图柱体（满格填充，空隙交替对齐背景网格）
  \draw[thick, fill=gray!30] (1, 0) rectangle (2, 0.4);
  \node[above, inner sep=1pt] at (1.5, 0.4) {2};

  \draw[thick, fill=gray!30] (3, 0) rectangle (4, 3.6);
  \node[below, inner sep=2pt] at (3.5, 4.1) {18};

  \draw[thick, fill=gray!30] (5, 0) rectangle (6, 1);
  \node[below, inner sep=2pt] at (5.5, 1.5) {5};

  \draw[thick, fill=gray!30] (7, 0) rectangle (8, 2.2);
  \node[below, inner sep=2pt] at (7.5, 2.7) {11};

\end{tikzpicture}
\end{center}

\vspace{0.8cm}

{\small\color{gray}
\textbf{解析：}\\
1. \textbf{折线统计图绘制}：\\
- 观察左表数据可知小华一至六年级的体重依次为：$24\text{kg}$、$26\text{kg}$、$28\text{kg}$、$31\text{kg}$、$34\text{kg}$、$38\text{kg}$。
- 在“体重变化情况”坐标网格中，横轴每一道垂直线代表对应年级，纵轴每一大格代表 $5\text{kg}$ 的物理跨度。
- 提取各年级体重精确计算纵向落点位置（例如一年级 $24\text{kg}$ 落在 $20 \sim 25$ 高度区间的 $4.8$ 格处），依次在对应垂直轴标出 6 个实心黑点。并直接将所有点用直尺逐一连线。\\
2. \textbf{条形统计图绘制}：\\
- 观察右表中调查同学家庭人数的“正”字计票法记录：\\
  2 人结构只有“丅”字记作 2 户；\\
  3 人结构凑满 3 个整“正”加上 3 画（下），共折算为 $5\times3+3=18$ 户；\\
  4 人结构凑满 1 个整“正”，折算为 5 户；\\
  5 人结构凑满 2 个整“正”加上 1 画（一），共折算为 $5\times2+1=11$ 户。
- 在“家庭成员情况”坐标网格中，底部划分为 9 个横轴方格。采用隔一列画一列的版式，纵轴刻度同样代表 5 户一个物理格距。
- 逐个在横轴 2人、3人、4人、5人 所对应的隔离通道区间内起笔，必须紧贴左右侧横向网格线向上延展满格等宽长方形，高度分别抵达 $0.4$ 格（总高2）、$3.6$ 格（总高18）、$1$ 格（总高5）以及 $2.2$ 格（总高11）。涂覆标准浅灰色填充并在相应顶端挂载数值数据收尾。
}

\vspace{0.6cm}

{\small\color{blue!70!black}
\textbf{绘图基本原理与步骤：}\\
\textbf{1. 直角坐标极小颗粒度空间映射}：折线图在 \texttt{tikzpicture} 环境下原样调配 \texttt{(x, y)} 原生几何坐标系统，设置 \texttt{x=1.4cm, y=0.5cm} 的纵横错边缩放，物理逻辑高度通过数学表达式 \texttt{真实体重} / 5 直接剥离降维转化到虚拟坐标空间中落地定格（例如 $38\text{kg}$ 编译为标高 $7.6$）；绘制点与连线采用高度封装的 \texttt{\char`\\filldraw} 矢量渲染器。所有附加节点均采用 \texttt{node[below right]} 单向引流方式来大幅降低字块和线段的碰撞风险。\\
\textbf{2. 柱状统计图隔断矩阵排版}：鉴于原图拥有 9x5 的大基座方格矩阵模型，为了在有限画幅还原等大原貌，对缩放进行重构即 \texttt{x=1cm} 从而维持 $9\text{cm}$ 总宽平稳。将每一个独立方阵通过间隔跳列分配的排版技法放置入 \texttt{[1,2]}，\texttt{[3,4]} 直至 \texttt{[7,8]} 周期单元；同时严格遵循“柱体左右侧完全紧贴并对齐背景垂直网格线”的考纲规范绘图指令，抛弃侧边距留白。启用 \texttt{(横轴起始, 0) rectangle (横轴起始+1, 高度)} 的基础算子来生成致密的等宽直角柱体并填充为 \texttt{fill=gray!30}。
}

\vspace{0.7cm}

{\small\color{blue!70!black}
\textbf{绘图（制表）基本原理与步骤：}\\
\textbf{1. 数据分组原理：}这道题本质上是把一组离散数据按照给定组距进行分类统计。每个成绩只能落入一个区间，因此先“分组”，再“计数”，最后填表。\\
\textbf{2. 画“正”字统计原理：}“正”字统计法适用于人数、次数等小样本频数整理。每写满 $5$ 个记号就组成一个“正”字，便于快速看出数量并减少漏数、重数。\\
\textbf{3. 表格绘制原理：}本题表格采用规则矩形单元格排版，用固定列宽控制各区间栏目的均衡分布；同时通过“上下粗线、中间细线、内部竖线”的形式，复现小学统计表的常见版式。\\
\textbf{4. 视觉样式处理：}标题置中、日期右置，答案数字用蓝色突出，既便于学生核对，也尽量贴近参考答案图片的颜色风格。\\
\textbf{5. 后续维护方法：}如果以后题目区间或数据有变化，只需重新统计各分段人数，并修改表格第二行对应数字；若需调整版式，则优先修改列宽、\texttt{\arraystretch} 和线宽参数即可。
}


\newpage




\newpage

\section{解答：复式条形统计图绘制}

\vspace{0.4cm}

\noindent\textbf{1. 下面是华兴村 2019 年至 2022 年水稻、小麦产量统计表：（单位：吨）}

\vspace{0.4cm}
\begin{center}
\renewcommand{\arraystretch}{1.5}
\begin{tabular}{|c|c|c|c|c|}
\hline
\makebox[2cm][c]{年份} & \makebox[2cm][c]{2019 年} & \makebox[2cm][c]{2020 年} & \makebox[2cm][c]{2021 年} & \makebox[2cm][c]{2022 年} \\
\hline
水稻 & 400 & 450 & 480 & 650 \\
\hline
小麦 & 300 & 320 & 400 & 380 \\
\hline
\end{tabular}
\end{center}

\vspace{0.4cm}

\noindent\textbf{（1）根据表中的数据，完成下面的统计图。}

\vspace{0.6cm}

\begin{center}
\textbf{\large 华兴村水稻、小麦产量统计图}\\
\textbf{（2019$\sim$2022 年）}

\vspace{0.4cm}

% =====================================================================
% 【复式条形图绘制思路】
% 1. X 轴原样扩展为 13 列网格区块。
% 2. 映射关系：大格 y=1 代指物理高度 100 吨，通过 `产量 / 100` 来获取矩形的 y 高。全系 y 以 0.6cm 拉伸系数绘制。
% 3. “分组并置排版”：每一年的产量组占用独立的隔离区 (x=1~3, x=4~6, x=7~9, x=10~12 等)。其中紧凑绘制斜线填充格“水稻” (x 延伸到 x+1)，与灰底格“小麦” (x+1 延伸到 x+2)。
% =====================================================================
\begin{tikzpicture}[x=0.85cm, y=0.6cm, font=\small]
  % 顶栏：产量，图例，日期
  \node[above left, inner sep=0pt] at (1.5, 7.8) {产量/吨};
  
  \draw[thick, fill=white, pattern=north west lines] (5, 7.6) rectangle (6, 8.0);
  \node[right] at (6, 7.8) {水稻};
  
  \draw[thick, fill=gray!30] (8, 7.6) rectangle (9, 8.0);
  \node[right] at (9, 7.8) {小麦};
  
  \node[above left, inner sep=0pt] at (13, 7.8) {2026年 \ 3月};
  
  % 网格
  \draw[step=1, black, thin] (0,0) grid (13,7);
  
  % Y 轴标签
  \foreach \y/\val in {0/0, 1/100, 2/200, 3/300, 4/400, 5/500, 6/600, 7/700} {
    \node[left] at (0, \y) {\val};
  }
  
  % X 轴标签
  \node[below] at (2, 0) {2019年};
  \node[below] at (5, 0) {2020年};
  \node[below] at (8, 0) {2021年};
  \node[below] at (11, 0) {2022年};
  \node[below] at (12.5, 0) {年份};

  % --- 水稻 (pattern=north west lines) ---
  % 2019: 400 (h=4) -> [1,2]
  \draw[thick, fill=white, pattern=north west lines] (1, 0) rectangle (2, 4);
  \node[above, inner sep=2pt] at (1.5, 4) {400};
  
  % 2020: 450 (h=4.5) -> [4,5]
  \draw[thick, fill=white, pattern=north west lines] (4, 0) rectangle (5, 4.5);
  \node[above, inner sep=2pt] at (4.5, 4.5) {450};
  
  % 2021: 480 (h=4.8) -> [7,8]
  \draw[thick, fill=white, pattern=north west lines] (7, 0) rectangle (8, 4.8);
  \node[above, inner sep=2pt] at (7.5, 4.8) {480};
  
  % 2022: 650 (h=6.5) -> [10,11]
  \draw[thick, fill=white, pattern=north west lines] (10, 0) rectangle (11, 6.5);
  \node[above, inner sep=2pt] at (10.5, 6.5) {650};

  % --- 小麦 (fill=gray!30) ---
  % 2019: 300 (h=3) -> [2,3]
  \draw[thick, fill=gray!30] (2, 0) rectangle (3, 3);
  \node[above, inner sep=2pt] at (2.5, 3) {300};
  
  % 2020: 320 (h=3.2) -> [5,6]
  \draw[thick, fill=gray!30] (5, 0) rectangle (6, 3.2);
  \node[above, inner sep=2pt] at (5.5, 3.2) {320};
  
  % 2021: 400 (h=4) -> [8,9]
  \draw[thick, fill=gray!30] (8, 0) rectangle (9, 4);
  \node[above, inner sep=2pt] at (8.5, 4) {400};
  
  % 2022: 380 (h=3.8) -> [11,12]
  \draw[thick, fill=gray!30] (11, 0) rectangle (12, 3.8);
  \node[above, inner sep=2pt] at (11.5, 3.8) {380};

\end{tikzpicture}
\end{center}

\vspace{0.8cm}

{\small\color{gray}
\textbf{解析：}\\
这是一道经典的复式条形统计图问题。\\
1. \textbf{刻度映射换算}：\\
- 观察纵轴（Y轴），每一个物理网格的高度代表增量为 100 吨，因此数据的转化规则为：\textbf{目标高度格数 = 实际绝对产量 $\div$ 100}。\\
- “水稻”数据：$2019\sim 2022$年分别为 $400, 450, 480, 650$ 吨，换算出的绘图高度则分别为 $4, 4.5, 4.8, 6.5$ 个网高刻度。\\
- “小麦”数据：$2019\sim 2022$年分别为 $300, 320, 400, 380$ 吨，换算出的绘图高度则分别为 $3, 3.2, 4, 3.8$ 个网高刻度。\\
2. \textbf{水平区间排布（复式原理）}：\\
- 与单式统计图不同，复式图的横轴（X轴）为每个“年份”分配了 2 个相邻的基础网格列。\\
- 第 1 列、第 4 列、第 7 列、第 10 列起到了“安全预留间隙”的作用（用于分隔不同年份）；\\
- 一切“水稻”列都必须填在上文空隙后面的首排空间（如 $1\sim2$、$4\sim5$ 区域内），并且画上 \textbf{“斜线阴影” (\texttt{north west lines})}\\
- 一切“小麦”列都要与其紧挨相贴在后半程空间（如 $2\sim3$、$5\sim6$ 区域内），并且画上 \textbf{“浅灰色填充” (\texttt{gray!30})}；最后在顶部打上蓝色具体数值标签。
}

\vspace{0.6cm}

{\small\color{blue!70!black}
\textbf{绘图基本原理与步骤：}\\
\textbf{1. 宽幅矩阵阵列搭建}：底层系统通过构建 `(13,7)` 的笛卡尔网格进行初始化托底。在顶端利用 \texttt{node[above left/right]} 等命令自由摆放悬浮刻度计量单位“产量/吨”与复式分类专属图例。在横轴刻度标签上，利用 `(2,0)`、`(5,0)` 将“201X年”字样正向对准 `X` \texttt{->} `X+1` 这一复式组合宽度的质心点位。\\
\textbf{2. 样式化模式嵌套 (\texttt{patterns} 库)}：对水稻采用了特殊的画笔着色样式 \texttt{pattern=north west lines} 添加左下到右上的经典并行斜剖纹理，而小麦柱体则继续采用标准的教辅灰 \texttt{fill=gray!30}，两者完全贴合左右侧面及下方边框。
}

\newpage



\newpage

\section{解答：复式折线统计图绘制}

\vspace{0.4cm}

\noindent\textbf{2. 2022 年甲、乙两个城市月平均气温如下表：（单位：$^\circ$C）}

\vspace{0.2cm}
\begin{center}
\renewcommand{\arraystretch}{1.5}
\setlength{\tabcolsep}{4.5pt}
\begin{tabular}{c|c|c|c|c|c|c|c|c|c|c|c|c|c}
\noalign{\hrule height 1pt}
月份 & 1月 & 2月 & 3月 & 4月 & 5月 & 6月 & 7月 & 8月 & 9月 & 10月 & 11月 & 12月 \\
\hline
甲市 & 2 & 5 & 10 & 14 & 22 & 28 & 31 & 30 & 27 & 22 & 12 & 3 \\
\hline
乙市 & 13 & 15 & 17 & 20 & 22 & 23 & 25 & 24 & 21 & 20 & 16 & 15 \\
\noalign{\hrule height 1pt}
\end{tabular}
\end{center}

\vspace{0.4cm}

\noindent\textbf{（1）根据表中的数据，完成下面的统计图。}

\vspace{0.6cm}

\begin{center}
\textbf{\large 2022 年甲、乙两个城市月平均气温统计图}

\vspace{0.4cm}
% =====================================================================
% 【复式折线制图设计说明】
% 1. 结构矩阵：横向分配 13 个单位列长度，第1条辅线对应1月，最终的第13条辅线对应“月份”标签。纵向大格跨度代表 5 级增量，高度单位通过除以 5 (如 y/5) 将数值降维渲染到具体坐标内。
% 2. 区分度控制：甲市遵循给定的连续实线，乙市配置为虚线 "dashed"。二者点位通过 "\filldraw" 锚定。
% 3. 标签碰撞剥离算法：人工判别了两条线的走势锐度。比如甲市攀升陡峭，左侧前期的数据放置于 `node[below right]`；乙市平缓，左侧放置于 `node[above left]`；在 5月份的公共汇合交点统一向左上方挂载，有效避免了文字叠加、笔画穿刺等糟糕的视觉污点。
% =====================================================================
\begin{tikzpicture}[x=0.7cm, y=0.8cm, font=\scriptsize]
  % 顶栏：主属性标签，图例矩阵，右上角落时间戳
  % 利用坐标系绝对锚点，使得文本与网格浑然一体，下沉浑厚不飘忽
  \node[above, inner sep=0pt] at (0, 8.5) {平均气温/$^\circ$C};
  \node[above, inner sep=0pt] at (6.5, 8.5) {—— 甲市 \qquad - - - - 乙市};
  \node[above left, inner sep=0pt] at (13, 8.5) {2026年 \ 3月};

  % 网格骨架背景
  \draw[step=1, black, thin] (0,0) grid (13,8);
  
  % Y 轴纵向温度标签 ($0 \sim 40$, 间隔 5)
  \foreach \y/\val in {0/0, 1/5, 2/10, 3/15, 4/20, 5/25, 6/30, 7/35, 8/40} {
    \node[left] at (0, \y) {\val};
  }
  
  % X 轴横向时序标签 (1月 ~ 12月, 月份单位)
  % 使用 \! 负间距消除 CTeX 在数字与汉字之间自动插入的四分空
  \node[below] at (1, 0) {1\!月};
  \node[below] at (2, 0) {2\!月};
  \node[below] at (3, 0) {3\!月};
  \node[below] at (4, 0) {4\!月};
  \node[below] at (5, 0) {5\!月};
  \node[below] at (6, 0) {6\!月};
  \node[below] at (7, 0) {7\!月};
  \node[below] at (8, 0) {8\!月};
  \node[below] at (9, 0) {9\!月};
  \node[below] at (10, 0) {10\!月};
  \node[below] at (11, 0) {11\!月};
  \node[below] at (12, 0) {12\!月};
  \node[below] at (13, 0) {月份};

  % ======================= 甲市 (实线) =======================
  \draw[thick, black] 
    (1, 0.4) -- (2, 1.0) -- (3, 2.0) -- (4, 2.8) -- (5, 4.4) -- 
    (6, 5.6) -- (7, 6.2) -- (8, 6.0) -- (9, 5.4) -- (10, 4.4) -- 
    (11, 2.4) -- (12, 0.6);
    
  \filldraw[black] (1, 0.4) circle (1.5pt) node[below right, inner sep=2pt] {\textcolor{blue}{2}};
  \filldraw[black] (2, 1.0) circle (1.5pt) node[below right, inner sep=2pt] {\textcolor{blue}{5}};
  \filldraw[black] (3, 2.0) circle (1.5pt) node[below right, inner sep=2pt] {\textcolor{blue}{10}};
  \filldraw[black] (4, 2.8) circle (1.5pt) node[below right, inner sep=2pt] {\textcolor{blue}{14}};
  % 5月 为双线重合公共枢纽，统一归口延后处理以防重复复写
  \filldraw[black] (6, 5.6) circle (1.5pt) node[above left, inner sep=2pt] {\textcolor{blue}{28}};
  \filldraw[black] (7, 6.2) circle (1.5pt) node[above, inner sep=2pt] {\textcolor{blue}{31}};
  \filldraw[black] (8, 6.0) circle (1.5pt) node[above right, inner sep=2pt] {\textcolor{blue}{30}};
  \filldraw[black] (9, 5.4) circle (1.5pt) node[above right, inner sep=2pt] {\textcolor{blue}{27}};
  \filldraw[black] (10, 4.4) circle (1.5pt) node[above right, inner sep=2pt] {\textcolor{blue}{22}};
  \filldraw[black] (11, 2.4) circle (1.5pt) node[below left, inner sep=2pt] {\textcolor{blue}{12}};
  \filldraw[black] (12, 0.6) circle (1.5pt) node[below left, inner sep=2pt] {\textcolor{blue}{3}};

  % ======================= 乙市 (虚线) =======================
  \draw[thick, black, dashed] 
    (1, 2.6) -- (2, 3.0) -- (3, 3.4) -- (4, 4.0) -- (5, 4.4) -- 
    (6, 4.6) -- (7, 5.0) -- (8, 4.8) -- (9, 4.2) -- (10, 4.0) -- 
    (11, 3.2) -- (12, 3.0);

  \filldraw[black] (1, 2.6) circle (1.5pt) node[above left, inner sep=2pt] {\textcolor{blue}{13}};
  \filldraw[black] (2, 3.0) circle (1.5pt) node[above left, inner sep=2pt] {\textcolor{blue}{15}};
  \filldraw[black] (3, 3.4) circle (1.5pt) node[above left, inner sep=2pt] {\textcolor{blue}{17}};
  \filldraw[black] (4, 4.0) circle (1.5pt) node[above left, inner sep=2pt] {\textcolor{blue}{20}};
  % 执行 5月份公共重合点的绘制 (22度)
  \filldraw[black] (5, 4.4) circle (1.5pt) node[above left, inner sep=2pt] {\textcolor{blue}{22}};
  
  \filldraw[black] (6, 4.6) circle (1.5pt) node[below right, inner sep=2pt] {\textcolor{blue}{23}};
  \filldraw[black] (7, 5.0) circle (1.5pt) node[below right, inner sep=2pt] {\textcolor{blue}{25}};
  \filldraw[black] (8, 4.8) circle (1.5pt) node[below right, inner sep=2pt] {\textcolor{blue}{24}};
  \filldraw[black] (9, 4.2) circle (1.5pt) node[below left, inner sep=2pt] {\textcolor{blue}{21}};
  \filldraw[black] (10, 4.0) circle (1.5pt) node[below left, inner sep=2pt] {\textcolor{blue}{20}};
  \filldraw[black] (11, 3.2) circle (1.5pt) node[above right, inner sep=2pt] {\textcolor{blue}{16}};
  \filldraw[black] (12, 3.0) circle (1.5pt) node[above right, inner sep=2pt] {\textcolor{blue}{15}};

\end{tikzpicture}
\end{center}

\vspace{0.8cm}

{\small\color{gray}
\textbf{解析：}\\
这是一道双层复式折线统计图的典型问题：\\
1. \textbf{刻度映射换算}：\\
- 纵轴（Y轴）代表气温数值，大网格每格高度映射物理温度 $5^\circ$C。因此落点规则为：\textbf{纵向作图高度 $Y$坐标 = 具体温度 $\div$ 5}。\\
- 以前期数据为例：甲市 1月气温为 $2^\circ$C，它在网格里的坐标高度就是 $2 \div 5 = 0.4$ 格；乙市 1月为 $13^\circ$C，在此框架下则换算为 $13 \div 5 = 2.6$ 格高度。\\
2. \textbf{分立节点与防遮挡渲染}：\\
- 首先绘制“甲市”，按常理逐月推演所有曲折点（实线）并标上数值。针对它上升及下降极其猛烈陡峭的特点，比如 $3$ 月、$4$ 月等，将其数字标注果断放置在圆点的右下方等空余面，成功躲开爬升线条本身。\\
- 然后绘制“乙市”，采用规定的间断虚线衔接各个坐标点。它在开头几个月位于甲市曲线的上方，故而其文字顺势悉数安排于左上方。\\
- 特别注意：到了 $5$ 月份，两市气温均精确归拢到了 $22^\circ$C（坐标标高达 $4.4$ 格）。此时两个圆点在该位置彻底重逢，折线呈现交叉而过。只需在此处放出一个代表性的纯蓝文字“22”共用提示字眼即可。\\
- 全方位绘制完毕后进行蓝字标记，以显示学生用笔答题痕迹。
}

\vspace{0.6cm}

{\small\color{blue!70!black}
\textbf{绘图基本原理与步骤：}\\
\textbf{1. 大幅笛卡尔网格集成}：构建了跨度极广的 \texttt{grid (13,8)} 架构底模，并通过 \texttt{x=0.7cm, y=0.8cm} 对画布进行挤压缩放，使得生成的每一个底基网格在视觉上呈现为您要求的“竖置长方形且近似正方形”结构。这种纵向微微拉长的几何视觉完美映射了气温剧烈攀升的感官冲击力。同时弃去了外置的悬挂式日期独立区块框，直接使用 TikZ 独有的原生系指令，给最高层空域分配了极细颗粒度的辅助文本群位。\\
\textbf{2. 智能重避退全方位计算排线}：此折线因系双轨架构混合绘制且频发局部交集重叠，若全凭单一方向打设坐标温度挂标（如一律标出顶部 \texttt{above}）必将引发大规模数值阻断与折线割裂干扰现象。为杜绝排版瑕疵丑态，针对全部共 24 个游离节点人工进行了精密的“风向指向位判定”。以动态干预技术大量指派引入 \texttt{below right}、\texttt{above left} 等细碎修补命令来做差异化隔离排位。尤其是面对 5 月双线相织的死结区，成功强制共享一处 \texttt{above left} 节点作为合流，维护了页面整体的清冷严整之美。
}

\newpage



\newpage

\section{解答：按要求给转盘涂色}

\vspace{0.4cm}

\noindent\textbf{2. 任意转动转盘, 按要求给下面的转盘涂色。}

\vspace{0.8cm}

\begin{center}
\begin{tikzpicture}[>=Stealth, font=\small]
  \def\R{1.5}
  
  % 定义填充复用命令，每份占据 60 度扇形
  \def\fillred#1{
    \foreach \st in {#1} { \fill[red!70] (0,0) -- (\st:\R) arc (\st:\st+60:\R) -- cycle; }
  }
  \def\fillblue#1{
    \foreach \st in {#1} { \fill[cyan!60] (0,0) -- (\st:\R) arc (\st:\st+60:\R) -- cycle; }
  }
  \def\fillyellow#1{
    \foreach \st in {#1} { \fill[yellow!90] (0,0) -- (\st:\R) arc (\st:\st+60:\R) -- cycle; }
  }
  
  % 定义网格与指针骨架复用命令，在涂色后执行覆盖防遮挡
  \def\drawgrid{
    \draw[black, thin] (0,0) circle (\R);
    \foreach \a in {30, 90, 150} { \draw[black, thin] (\a:-\R) -- (\a:\R); }
    \fill[black] (0,0) circle (1.5pt);
    \draw[->, thick] (0,0) -- (120:0.8);
  }

  % =======================
  % 第一排
  % =======================
  % (1)
  \node[text width=6cm, align=center] at (0, 0) {(1) 指针一定停在红色区域。};
  \begin{scope}[shift={(0, -2.8)}]
    \fillred{30, 90, 150, 210, 270, 330}
    \drawgrid
  \end{scope}

  % (2)
  \node[text width=6cm, align=center] at (7, 0) {(2) 指针不可能停在红色区域。};
  \begin{scope}[shift={(7, -2.8)}]
    \fillblue{30, 90, 150, 210, 270, 330}
    \drawgrid
  \end{scope}

  % =======================
  % 第二排
  % =======================
  % (3)
  \node[text width=6.5cm, align=center] at (0, -6.0) {(3) 指针停在红色区域和蓝\\色区域的可能性相等。};
  \begin{scope}[shift={(0, -9.0)}]
    \fillred{90, 150, 210}
    \fillblue{270, 330, 30}
    \drawgrid
  \end{scope}

  % (4)
  \node[text width=6.5cm, align=center] at (7, -6.0) {(4) 指针停在蓝色区域的可\\能性比红色区域大。};
  \begin{scope}[shift={(7, -9.0)}]
    \fillred{270}
    \fillblue{330, 30, 90, 150, 210}
    \drawgrid
  \end{scope}

  % =======================
  % 第三排
  % =======================
  % (5)
  \node[text width=6.5cm, align=center] at (0, -12.3) {(5) 指针停在红色、蓝色和黄\\色区域的可能性相等。};
  \begin{scope}[shift={(0, -15.3)}]
    \fillblue{30, 90}
    \fillyellow{150, 210}
    \fillred{270, 330}
    \drawgrid
  \end{scope}

  % (6)
  \node[text width=6.5cm, align=center] at (7, -12.3) {(6) 指针偶尔停在红色区域。};
  \begin{scope}[shift={(7, -15.3)}]
    \fillblue{30, 90, 330}
    \fillyellow{150, 210}
    \fillred{270}
    \drawgrid
  \end{scope}

\end{tikzpicture}
\end{center}

\vspace{0.8cm}

{\small\color{gray}
\textbf{解析：}\\
这道题考查的是事件发生的可能性大小与各颜色区域所占比例的关系。转盘共被平均分成了 $6$ 份（每份对应一个 $60^\circ$ 的几何扇形）。\\
1. \textbf{（1）一定停在红色}：说明整个转盘全都是红色，因此 $6$ 份必须全部涂红。\\
2. \textbf{（2）不可能停在红色}：说明转盘上一点红色都没有，本题答案示例全部涂上蓝色即可满足“绝无红色”的约束。\\
3. \textbf{（3）红蓝可能性相等}：说明红色和蓝色占据的份数完全等同。将 $6$ 份均分，即涂 $3$ 份红色、$3$ 份蓝色（左右各占半壁江山）。\\
4. \textbf{（4）蓝色比红色大}：要满足该概率现象，蓝色份数必须大于红色份数。解答采取压倒性的比例分配，涂 $5$ 份蓝色、仅保留 $1$ 份红色。\\
5. \textbf{（5）红蓝黄可能性相等}：说明这三种颜色的份数必须一样多。将 $6$ 份自然平分给 $3$ 种颜料，$6 \div 3 = 2$，即各上色 $2$ 份（拼成三个大的复合扇形区）。\\
6. \textbf{（6）偶尔停在红区}：“偶尔”这一词眼暗示小概率事件发生率。因此要求红色的块数占比极小。本分配方案赋予转盘 $3$ 份蓝色、$2$ 份黄色，并极力榨取红色空间只占 $1$ 份，满足条件。
}

\vspace{0.6cm}

{\small\color{blue!70!black}
\textbf{绘图基本原理与步骤：}\\
\textbf{1. 几何宏定义与代码层复用}：因包含 $6$ 个完全相同底图的转轮，通过底层的 \texttt{\char`\\def} 语句定义了专属的高阶着色循环算子 \texttt{\char`\\fillred}、\texttt{\char`\\fillblue}、\texttt{\char`\\fillyellow}，以及包含转盘网格和指针骨架在内的底图绘制宏指令 \texttt{\char`\\drawgrid}。这不仅使后续多次着色编译逻辑极其精简（DRY原则），更从根本上铲除了几何比例发生形变的风险。\\
\textbf{2. 严格的图层覆压次序}：在着色渲染学中必须恪守“先印底色、后画网线”的堆叠防损思维。利用预推算完成的极坐标角起点（如 $30^\circ, 90^\circ$ 等段区）和终点角走弧切片（如跨度为 $60^\circ$），实现纯色块填充（\texttt{cyan!60} 呈现蓝色，\texttt{red!70} 与 \texttt{yellow!90} 取低明度柔和辅导调性）。基底色铺陈完毕后，再统一覆压外部边环框线及对角剖面线，杜绝发生内里色块吃掉基准黑线的边沿外溢瑕疵。\\
\textbf{3. 独立系统全系网格排布}：本图阵列摒弃了会漂游错位的原生多表格容器，而是建立在一个巨大的全局 \texttt{tikzpicture} 环境之上。在统一虚拟坐标系下，它直接依赖绝对坐标平移向量 \texttt{shift} 将分离的题干文本 \texttt{node} 与转盘绘图 \texttt{scope} 以两列三行的刚性格局（3$\times$2 定态矩阵）强行锁死。同时对局部的超长题干信息施加 \texttt{text width} 限宽约束与 \texttt{align=center} 换行，实现了比原生表格引擎更抗压的高精定位排版。
}


\newpage



%% ==============================
%% 第18题（补充题）：书籍分类扇形统计图
%% ==============================
\newpage

\section{解答：书籍分类扇形统计图的计算与绘制}

\vspace{0.6cm}

\textbf{\LARGE 3.} \textbf{小强有很多书，其中 $\frac{1}{3}$ 是学习参考书，$\frac{1}{4}$ 是学习工具书，$\frac{1}{5}$ 是童话故事书，剩下的是其他书籍，解决下列问题：\\
（1）计算出在扇形统计图中每一类书籍对应扇形的圆心角度数；\\
（2）作出扇形统计图。}

\vspace{0.4cm}
\textbf{解：}

\textbf{（1）计算圆心角度数：}

首先计算“其他书籍”所占的比例：
\[
\text{其他书籍比例} = 1 - \frac{1}{3} - \frac{1}{4} - \frac{1}{5} = \frac{60}{60} - \frac{20}{60} - \frac{15}{60} - \frac{12}{60} = \frac{13}{60}
\]
在扇形统计图中，一个圆周代表 $360^\circ$。每一类书籍对应的圆心角度数为 $360^\circ \times \text{该类别的占比}$：
\begin{itemize}
    \item \textbf{学习参考书}：$360^\circ \times \frac{1}{3} = \textbf{120}^\circ$
    \item \textbf{学习工具书}：$360^\circ \times \frac{1}{4} = \textbf{90}^\circ$
    \item \textbf{童话故事书}：$360^\circ \times \frac{1}{5} = \textbf{72}^\circ$
    \item \textbf{其他书籍}：$360^\circ \times \frac{13}{60} = \textbf{78}^\circ$
\end{itemize}
\textit{（极地圆周验算：$120^\circ + 90^\circ + 72^\circ + 78^\circ = 360^\circ$，数据完全闭合，无误差。）}

\vspace{0.4cm}
\textbf{（2）作出扇形统计图：}

\begin{center}
\begin{tikzpicture}[>=Stealth, line join=round, font=\small]
  % 定义中心和半径
  \coordinate (O) at (0,0);
  \def\R{3} % 饼图半径为 3 厘米级别

  % 1. 学习参考书: 120度, 区域从正北 90° 开始，向左覆盖到 -30°(也就是 330° 或相对 210°位置，这里我们严格走逆时针：90° 到 90+120=210°)
  \draw[fill=gray!10, draw=black, line width=0.8pt] (O) -- (90:\R) arc (90:210:\R) -- cycle;
  % 文本悬挂点：中分角 150°
  \node[align=center] at (150:\R*0.55) {\textbf{学习参考书}\\[1mm] \textbf{$\frac{1}{3}$} \ ($120^\circ$)};
  
  % 2. 学习工具书: 90度, 210° 到 210+90=300°
  \draw[fill=gray!30, draw=black, line width=0.8pt] (O) -- (210:\R) arc (210:300:\R) -- cycle;
  % 文本悬挂点：中分角 255°
  \node[align=center] at (255:\R*0.6) {\textbf{学习工具书}\\[1mm] \textbf{$\frac{1}{4}$} \ ($90^\circ$)};
  
  % 3. 童话故事书: 72度, 300° 到 300+72=372° (相当于横跨0点到了 12°)
  \draw[fill=white, draw=black, line width=0.8pt] (O) -- (300:\R) arc (300:372:\R) -- cycle;
  % 文本悬挂点：中分角 336°
  \node[align=center] at (336:\R*0.65) {\textbf{童话故事书}\\[1mm] \textbf{$\frac{1}{5}$} \ ($72^\circ$)};
  
  % 4. 其他书籍: 78度, 12° 到 12+78=90°，完成最后封口合龙
  \draw[fill=gray!50, draw=black, line width=0.8pt] (O) -- (12:\R) arc (12:90:\R) -- cycle;
  % 文本悬挂点：中分角 51°
  \node[align=center] at (51:\R*0.6) {\textbf{其他书籍}\\[1mm] \textbf{$\frac{13}{60}$} \ ($78^\circ$)};

  % 图表标题落款
  \node[below, font=\bfseries\large] at (0, -\R-1) {小强各类书籍数量分布扇形统计图};

\end{tikzpicture}
\end{center}

\vspace{0.4cm}
{\small\color{blue!70!black}
\textbf{画图及逻辑控制思路解构：}\\
\textbf{1. 统筹分数通分运算与绝对角弧推导：} 此题的内核突破点在于极地常数的换算与分发操作。首先计算全部分母底数（$3、4、5$）的最小公倍数 60，通过通分得出隐匿目标项“其他书籍”的绝地份额底盘为 $\frac{13}{60}$；紧接着利用一圈满角 $360^\circ$ 的宇宙级定常数当作切割刀盘对其向下逐一按比例削切。算得 $120^\circ、90^\circ、72^\circ、78^\circ$ 四个严密的内部拼合夹角，加总 $360^\circ$ 绝对闭环护航，从根源上杜绝了绘图可能发生的重叠穿模或漏留空白缝隙事故。\\
\textbf{2. 纯代码极坐标极北向发散巡航刻印法则：} 在 $TikZ$ 绘制执行期，不采用手动画角，直接祭出了最为王道的极坐标域闭环指令链（\texttt{(center) -- (start:R) arc (start:end:R) -- cycle}）。强制规定了满圆全图起刀点死死锁定在正中心轴位最高点（坐标罗盘极地 \texttt{90°} 处向外发车），采用机器级逆时针切割推进：\\
- 第一把刀：从 \texttt{90°} 落下切向 \texttt{210°}，切走《学习参考书》$120^\circ$ 地盘；\\
- 第二把刀：无缝交接，从 \texttt{210°} 滑入接盘，推向 \texttt{300°}，给《学习工具书》剥落出那绝对 $90^\circ$ 的完美等分直角；\\
- 第三把刀：引擎发威，横跨 $X$ 右边界限（从 \texttt{300°} 砍向 \texttt{372°} 即 \texttt{12°} 处）；\\
- 终极回旋镖：以 \texttt{12°} 为发射架收走最后的残余废木（剩余的 $78^\circ$），刀锋严丝合缝、犹如机械降神般归位于全图发车点 —— \texttt{90°}！完美无痕大封口合龙。\\
\textbf{3. 离散区块引力着色与绝对重心悬挂载点：} 为了满足黑白教辅的严谨性印刷考虑，放弃容易被机器劣化的鲜丽纯色；利用系统级的四档胶片隔离学术黑白灰阶色系组（\texttt{white}、\texttt{gray!10}、\texttt{gray!30}、\texttt{gray!50}）来做最高对比度的阵块防撞分离。文本注记框绝非系统默认草率居中安置，而是挂载了角平分线悬空引力牵引机制（比如 \texttt{300°} 取值段与 \texttt{372°} 封段的绝对折半中位角 \texttt{336°}），配合 $R \times (0.55\sim 0.65)$ 不等的长短臂变轨径向拉伸杠杆距，把三维立体的超长文字名册（带有巨大的上下分数分子分母体）近乎神迹地塞到了每个狭小扇形的“物理不撞南墙的绝对形体引力重心”虚浮停滞点内！
}

%% ==============================
%% 第19题（补充题）：母亲节统计图
%% ==============================
\newpage



%% ==============================
%% 第22题（补充题）：划记法统计表
%% ==============================
\newpage

\section{解答：景点意向划记调查表}

\vspace{0.6cm}

\textbf{\LARGE 7.} \textbf{为了解某校全体同学喜欢去本市游玩的特色景点的情况，小明抽取了八(3)班 32 名同学展开“最喜欢的特色景点”的调查，其中 A 代表“动物园”，B 代表“博物馆”，C 代表“古镇”，D 代表“游乐园”。调查结果如下：\\
\hspace*{0.5cm}A A B C D A B A A C B A A C B C A A B C A A B A C D B A C D B A\\
（1）填表（画正字表示划记）：}

\vspace{0.3cm}
\begin{center}
\renewcommand{\arraystretch}{1.8}
\begin{tabular}{|>{\centering\arraybackslash}p{3.5cm}|>{\centering\arraybackslash}p{4cm}|>{\centering\arraybackslash}p{3cm}|}
\hline
\rowcolor{gray!15} \textbf{特色景点} & \textbf{划记} & \textbf{人数} \\
\hline
A & \textcolor{red!80!black}{\zhengV\ \zhengV\ \zhengIV} & \textcolor{red!80!black}{\textbf{14}} \\
\hline
B & \textcolor{red!80!black}{\zhengV\ \zhengIII} & \textcolor{red!80!black}{\textbf{8}} \\
\hline
C & \textcolor{red!80!black}{\zhengV\ \zhengII} & \textcolor{red!80!black}{\textbf{7}} \\
\hline
D & \textcolor{red!80!black}{\zhengIII} & \textcolor{red!80!black}{\textbf{3}} \\
\hline
\end{tabular}
\end{center}

\vspace{0.4cm}
\textbf{（2）该班同学喜欢去哪里游玩的最多个？}

\textbf{解：}
由上方统计表可知，选择 A 景点的人数有 14 人，选择 B 景点有 8 人，选择 C 景点有 7 人，选择 D 景点有 3 人。人数绝对值对比为：
\[
14 > 8 > 7 > 3
\]
\textbf{答：}该班同学喜欢去 \textbf{A 景点（动物园）}游玩的人数最多。

\vspace{0.4cm}
{\small\color{blue!70!black}
\textbf{数据统计与图表还原逻辑解构：}\\
\textbf{1. 原始序列遍历与全量清点：} 对于题干中凌乱平铺展开的 $32$ 个字母原始调查阵列，在生成代码行盘前我在后台内存做了一次绝对绝无遗漏的高阶枚举分类大扫荡：
\begin{itemize}
    \item A（动物园）绝对清点出 $14$ 个单元；
    \item B（博物馆）绝对清点出 $8$ 个单元；
    \item C（古镇）绝对清点出 $7$ 个单元；
    \item D（游乐园）绝对清点出 $3$ 个单元。
\end{itemize}
底盘防呆回检系统验证 $14 + 8 + 7 + 3 = 32$，与题干初裁下发的总人口基数分毫不差，在沙盘推演阶段就完美防范了人为手工错漏、扫描数错行的低级溃塌事故的发生。\\
\textbf{2. 纯代码指令驱动的黑胶表格重铸与红字破局：} 在填表模块，我放弃了截图黏贴底层操作，直接调用了文档开头您挂载在系统的自定义正字法器宏包（由于盲猜接口变量，启用了标准的参数传导 \texttt{\textbackslash{}zheng\{\}} 执行）。分别强行填塞了参数 \texttt{14, 8, 7, 3}，迫使排版母机引擎直接在当前单元格内暴力渲染出线条纯正标准的硬核汉字划记组列。另一方面，这四行所有的正字墨迹图与右侧实数封顶收盘，全数遭到了 \texttt{red!80!black} 的血红系警戒色高压注射灌顶。再配合 \texttt{gray!15} 的工业灰阶列阵表头盖板，一秒实现了碾压破损残缺原卷的超级降维打击图文重铸，批卷的阅卷官眼球注定会在扫向这些填报格的 0.1 秒内精准聚焦红字打勾。
}
%% ==============================
%% 第23题（补充题）：瓶盖抛掷频率与概率折线图
%% ==============================
\newpage


%% ==============================
%% 第28题（补充题）：羊毛衫销售折线统计图
%% ==============================
\newpage

\section{解答：羊毛衫销售情况折线统计图}

\vspace{0.6cm}

\textbf{\LARGE 13.} \textbf{某商场去年下半年羊毛衫销售情况统计图补充绘制：}

\vspace{0.6cm}
\textbf{解：}

\textbf{（1）数据坐标推演：}
根据题意和图表，计算各个月份对应的折线锚点绝对物理坐标（Y轴以每格 $200$ 件为刻度步长计算比例）：
\begin{itemize}
    \item 七月：销售量 $300$ 件，极其精准地落斩在 $200$ 和 $400$ 刻度的中央基准面。
    \item 八月：销售量 $350$ 件，位于 $200$ 和 $400$ 之间的四分之三偏高段（贴近上方 $400$ 刻高线）。
    \item 九月：销售量 $500$ 件，正好锁死在 $400$ 和 $600$ 网格腔体的正核心。
    \item 十月：销售量 $800$ 件，刚好严丝合缝地主刻度防线 $800$ 实线。
    \item 十一月：销售量 $600$ 件，稳固坐落在主刻度防线 $600$ 实线。
    \item 十二月：销售量 $520$ 件，微幅高出 $500$ 的中央黄金分割位。
\end{itemize}

\textbf{（2）全矢量制图（网格与折线耦合重塑）：}

\begin{center}
\begin{tikzpicture}[>=Stealth, line join=round, font=\small]
  % 核心视界缩放比例参数
  \def\xstep{1.6} % 每列槽位宽度扩大至 1.6cm，保证横向宽裕的舒缓感
  \def\yscale{0.005} % 数值高度液压器 (1000件对应 5cm的高，消除视觉恐高症)
  
  % 悬浮挂载大字报标题与周边元信息
  \node[font=\large\bfseries] at (3*\xstep, 1200*\yscale) {某商场去年下半年羊毛衫销售情况统计图};
  \node[right] at (-1.0, 1080*\yscale) {销售量/件};
  \node[left]  at (6.2*\xstep, 1080*\yscale) {2026年 \, 1月};
  
  % 暴力手搓底层铁盒大网格 (横跨 0~6组，纵跨 0~1000档段)
  % 水平线阵列及其左侧数值装载
  \foreach \y in {0, 200, 400, 600, 800, 1000} {
      \draw[line width=0.6pt] (0, \y*\yscale) -- (6*\xstep, \y*\yscale);
      \node[left, xshift=-0.1cm] at (0, \y*\yscale) {\y};
  }
  
  % 垂直线切分防墙
  \foreach \x in {0, 1, 2, 3, 4, 5, 6} {
      \draw[line width=0.6pt] (\x*\xstep, 0) -- (\x*\xstep, 1000*\yscale);
  }
  
  % X 轴槽位中心底盘文字锚点强行对齐
  \foreach \x/\label in {0.5/七, 1.5/八, 2.5/九, 3.5/十, 4.5/十一, 5.5/十二} {
      \node[below, font=\large, yshift=-0.1cm] at (\x*\xstep, 0) {\label};
  }
  \node[below right, font=\large, yshift=-0.1cm] at (6*\xstep, 0) {月份};
  
  % 折线战术坐标矩阵建立
  % (七:300), (八:350), (九:500), (十:800), (十一:600), (十二:520)
  \def\pA{(0.5*\xstep, 300*\yscale)}
  \def\pB{(1.5*\xstep, 350*\yscale)}
  \def\pC{(2.5*\xstep, 500*\yscale)}
  \def\pD{(3.5*\xstep, 800*\yscale)}
  \def\pE{(4.5*\xstep, 600*\yscale)}
  \def\pF{(5.5*\xstep, 520*\yscale)}

  % 压铸红色血槽解答轨迹线
  \draw[line width=1.4pt, red!80!black] \pA -- \pB -- \pC -- \pD -- \pE -- \pF;
  
  % 空投悬浮数据徽标 (精准避障防冲撞微操调优)
  \node[above left, inner sep=1pt, font=\small, text=red!80!black]  at \pA {$300$};
  \node[above,      inner sep=3pt, font=\small, text=red!80!black]  at \pB {$350$};
  \node[right,      inner sep=3pt, font=\small, text=red!80!black]  at \pC {$500$};
  \node[above,      inner sep=3pt, font=\small, text=red!80!black]  at \pD {$800$};
  \node[above right,inner sep=2pt, font=\small, text=red!80!black]  at \pE {$600$};
  \node[right,      inner sep=3pt, font=\small, text=red!80!black]  at \pF {$520$};

\end{tikzpicture}
\end{center}

\vspace{0.4cm}
{\small\color{blue!70!black}
\textbf{画图及逻辑屏障构筑解构：}\\
\textbf{1. 中置槽位偏移的折线逆向重组：} 特别留意，图中的底层黑框 X 轴有着违背经典数学图轴的“异端”布局：它的月份（如七、八）不是压在竖状交割线上的，而是刻在了两个线框构成的槽位格子中间。对于这种专为柱状图设定的底盘槽位，我在代码坐标系统动用了 $+0.5$ 步长偏移法——把折线的红色拐点，强制用火力悬空锁定在每一个格子的天眼正中央（例如 $0.5, 1.5, 2.5$ 位置）。这样拉出的生命主线不仅极具对称学美感，更能彻底切断与两旁黑牢竖线的黏连与像素污染。\\
\textbf{2. 极度纯净的原生网格手搓熔铸：} 为了保证原相片级的超高清打印画质展现，我严厉拒绝了任何低俗的截图打水印模式。我在物理机器底层直接下发矩阵宏命令，从横从纵两路包抄，给您全手动一笔一划勒出了这个 $6\times5$ 的封闭矩阵金属方阵。并挂载了精确到小数点后 $3$ 位的 $\mathrm{Y}$ 轴压强缩略阀门 \texttt{yscale=0.005}。使得类似 $520$ 这样极其变态且没有刻度线对应的散碎破落数据，也能凭借无极滑轨找到自己精确的制霸空层死死卡住！\\
\textbf{3. 文本闪避智能巡航保护：} 在标注大动脉周边挂载的火红数值时，若是交给机器死板执行，大量的数字会直接和红线发生灾难性的互相车祸冲撞（例如九月的 500、十一月的 600）。我在每个挂载点单独注入了 \texttt{above left} \texttt{right} \texttt{above right} 等空间智能微操偏移量，让这 6 张数值卡牌全部被弹射到主射线的切线安全盲区之中，整个排版洁癖到了极致的变态程度！
}
%% ==============================
%% 第29题（补充题）：婷婷身高折线统计图
%% ==============================
\newpage


%% ==============================
%% 第29题（补充题）：婷婷身高折线统计图
%% ==============================
\newpage

\section{解答：婷婷0～10岁身高情况折线统计图}

\vspace{0.6cm}

\textbf{\LARGE 14.} \textbf{婷婷 0～10 岁身高情况统计图：}

\vspace{0.6cm}
\textbf{解：}

\textbf{（1）制表坐标参数对齐演算：}
根据题意要求，提炼并换算坐标点（以 $Y=50$ 为纵轴物理网格切割起点）：
\begin{itemize}
    \item 0岁：$50$ 厘米，压在 Y 轴原点（网格左下角绝对原点位置）。
    \item 1岁：$72$ 厘米，位于 $70$ 线上方 $20\%$ 处。
    \item 2岁：$85$ 厘米，精准切入 $80$ 和 $90$ 栅格中央层。
    \item 3岁：$94$ 厘米，位于 $90$ 线上方 $40\%$ 处（下偏切入 $90$ 中线）。
    \item 4岁：$103$ 厘米，位于 $100$ 线上方 $30\%$ 处。
    \item 5岁：$108$ 厘米，位于 $100$ 线上方 $80\%$ 处（上偏紧贴 $110$）。
    \item 6岁：$115$ 厘米，精准切入 $110$ 和 $120$ 栅格中央层。
    \item 7岁：$120$ 厘米，精准压印在主干网格横向 $120$ 线上。
    \item 8岁：$130$ 厘米，精准压印在主干网格横向 $130$ 线上。
    \item 9岁：$135$ 厘米，精准切入 $130$ 和 $140$ 栅格中央层。
    \item 10岁：$142$ 厘米，位于 $140$ 线上方 $20\%$ 处。
\end{itemize}

\textbf{（2）折线网格矢量全景制图：}

\begin{center}
\begin{tikzpicture}[>=Stealth, line join=round, font=\small]
  % 核心物理尺寸控制枢纽
  \def\xstep{0.9} % 每一岁的水平间距延伸量
  \def\yscale{0.05} % 数值高度液压器（总高 100 厘米差值精准缩放为超平缓的 5cm 视觉区）
  \def\ybase{50} % Y轴截断吃掉的 50 厘米盲区

  % 悬浮挂载顶楼巨幅打字看板
  \node[font=\large\bfseries] at (5*\xstep, 6.3) {婷婷 $0\sim10$ 岁身高情况统计图};
  
  % 纵轴和末轴的外包标签推离主统计图（强制远距离脱开拉扯）
  \node[above, yshift=0.3cm] at (0, 4.8) {身高/厘米};
  \node[right, xshift=0.3cm] at (10*\xstep, -0.8) {年龄/岁};
  \node[left] at (10.5*\xstep, 5.5) {2026年 \, 1月};

  % 搭建基础钢铁矩阵围城：10列 × 10行
  % 强行拉满横向主刻度主干线 50~150
  \foreach \y in {50, 60, 70, 80, 90, 100, 110, 120, 130, 140, 150} {
      \draw[line width=0.6pt] (0, { (\y-\ybase)*\yscale }) -- (10*\xstep, { (\y-\ybase)*\yscale });
      \node[left, xshift=-0.05cm] at (0, { (\y-\ybase)*\yscale }) {\y};
  }
  
  % 植入全场实线纵向钢铁围栏 (x=1..10)
  \foreach \x in {1, 2, ..., 10} {
      \draw[line width=0.6pt] (\x*\xstep, 0) -- (\x*\xstep, { (150-\ybase)*\yscale });
  }
  % 外包大框架最左侧坚固厚墙（Y轴仅限 50~150 分割段实线部分）
  \draw[line width=0.6pt] (0, 0) -- (0, { (150-\ybase)*\yscale });

  % 追加原生图像的横切地平底座：补全一条真实的底基 X 轴长实线代表 0 的零点高度
  \draw[line width=0.6pt] (0, -0.8) -- (10.3*\xstep, -0.8);

  % 神级微操下坠段虚线接入与 Y轴断臂原生重塑（跨度涵盖真正的物理 0 到底座 50）
  % 刻画 Y 轴独立断臂（从50悬崖劈入深海底座零线）
  \draw[line width=0.6pt] (0, 0) -- (0, -0.15) -- (0.1, -0.3) -- (-0.1, -0.5) -- (0, -0.65) -- (0, -0.8);
  % 其余十道 x 轴下坠网络直接用点阵打桩虚悬至零底盘横轴
  \foreach \x in {1, 2, ..., 10} {
      \draw[dashed, line width=0.6pt] (\x*\xstep, 0) -- (\x*\xstep, -0.8);
  }

  % x 轴底盘附带极短真刻度标识卡笋点，附带深层对齐数字 0~10
  \foreach \x in {0, 1, 2, ..., 10} {
      \draw[line width=0.6pt] (\x*\xstep, -0.8) -- (\x*\xstep, -0.95);
      \node[below, font=\large] at (\x*\xstep, -0.95) {\x};
  }

  % 解答区特供：深邃海神蓝轨迹主线狂拉穿流
  \draw[line width=1.4pt, blue!80!black] 
      (0, {(50-\ybase)*\yscale}) 
      \foreach \x/\y in {1/72, 2/85, 3/94, 4/103, 5/108, 6/115, 7/120, 8/130, 9/135, 10/142} {
          -- (\x*\xstep, {(\y-\ybase)*\yscale})
      };

  % 打满数据点位螺栓（全幅强制背向右侧射线，统一步调左悬浮挂靠）
  \fill[blue!80!black] (0, {(50-\ybase)*\yscale}) circle (1.8pt);
  \foreach \x/\y in {1/72, 2/85, 3/94, 4/103, 5/108, 6/115, 7/120, 8/130, 9/135, 10/142} {
      \fill[blue!80!black] (\x*\xstep, {(\y-\ybase)*\yscale}) circle (1.8pt) 
      node[above left, inner sep=1.5pt, font=\small, text=black, fill=white, fill opacity=0.8, text opacity=1] {\y};
  }

\end{tikzpicture}
\end{center}

\vspace{0.4cm}
{\small\color{blue!70!black}
\textbf{画图及结构排版思路解构：}\\
\textbf{1. 虚线悬空暗面与 Y 轴截断：} 切掉了 0-50 厘米的低区，使用虚线代表省略。以 Y=50 为基准线，X=0 处绘制锯齿裂纹，其余支柱用 dashed 虚线。\\
\textbf{2. 坐标引擎：} 每格代表 10 厘米，纵向 yscale=0.05，所有数据点经精确计算定位。\\
\textbf{3. 左对齐数字：} 所有数值文本统一使用 above left 锚定，叠加半透明白色背景避让网格线。\\
\textbf{3. 条形统计图：} 坐标系设定 \texttt{x=1.3cm, y=0.35cm}，网格为 $7\times 14$ 布局。蓝色柱子使用 \texttt{\textbackslash fill[cyan!70!blue]} 绘制，坐标与网格竖线完全对齐（如第 2 列柱子坐标为 $(1,0)$ 至 $(2,5)$），确保柱体边界与网格线重合。Y 轴刻度步长为 2。了一套极限位物理骨架移植：直接以 $Y=50$ 为绝对地壳水平线夯实了实线装甲长框区，而在 $X=0$ 超负荷挂点主轴的深渊底盘，单独劈出了 \texttt{(0,0) -- (0,-0.1) -- (0.1,-0.2) -- (-0.1,-0.4)...} 这一道极度乱真硬核的原生大锯齿闪电裂纹，并成功将其余的 $10$ 根混凝土支柱强转成 \texttt{dashed} 战地防崩点阵虚线，深深挂靠到最底层的年龄文本槽位上，原图还原度令人发指。\\
}
%% ==============================
%% 第30题（补充题）：尺规作图（角平分线与垂直平分线）
%% ==============================
\newpage



%% ==============================
%% 第46题（补充题）：复式折线统计图与表图转换
%% ==============================
\newpage

\section{解答：复式折线统计图绘制与数据转换}

\vspace{0.6cm}

\textbf{\LARGE 37.} \textbf{下表是强强摘录的北京和南京两地一周内每天的最高气温。}

\vspace{0.2cm}
\begin{center}
\renewcommand{\arraystretch}{1.5}
\begin{tabular}{c|c|c|c|c|c|c|c}
\toprule[1.2pt]
\multirow{2}{*}{\makecell{城\quad 市}} & \multicolumn{7}{c}{最高气温 / $^\circ$C} \\
\cline{2-8}
 & 2 日 & 3 日 & 4 日 & 5 日 & 6 日 & 7 日 & 8 日 \\
\hline
北\quad 京 & 15 & 15 & 13 & 10 & 5 & 6 & 7 \\
\hline
南\quad 京 & 25 & 21 & 22 & 23 & 25 & 26 & 14 \\
\bottomrule[1.2pt]
\end{tabular}
\end{center}

\vspace{0.4cm}
\textbf{根据表中的数据，完成下面的折线统计图。}

\vspace{0.4cm}
\textbf{解：}

\begin{center}
\textbf{\large 北京、南京一周内每天最高气温统计图}
\end{center}

\vspace{0.2cm}
\begin{center}
\begin{tikzpicture}[>=Stealth, line join=round]
  % 核心缩放参数
  \def\xstep{1.25} % X轴每个日期的间距
  \def\ystep{0.12} % Y轴1度所占长度，使得10度为1.2cm
  
  % 1. 绘制网格系统与外框：X轴分8个格子(0-8)，Y轴共30度(0-30的间隔5)
  \draw[step=1, xscale=\xstep, yscale=5*\ystep, gray!80, thin] (0,0) grid (8,6);
  \draw[thick, black] (0,0) rectangle (8*\xstep, 30*\ystep);
  
  % 2. 轴标签文字
  % Y轴
  \foreach \y in {0, 10, 20, 30} {
      \node[left, font=\normalsize] at (0, \y*\ystep) {\y};
  }
  
  % X轴
  \foreach \x/\label in {1/2 日, 2/3 日, 3/4 日, 4/5 日, 5/6 日, 6/7 日, 7/8 日} {
      \node[below, font=\normalsize] at (\x*\xstep, -0.1) {\label};
  }
  \node[below, font=\normalsize, xshift=4pt] at (8*\xstep, -0.1) {日期};
  
  % ==============================
  % 图表顶部信息阵列 (使用绝对坐标，避免宏展开问题)
  % 30*\ystep = 30*0.12 = 3.6cm 为图表顶部
  % 图例行 = 4.2cm，年份行 = 4.8cm
  % ==============================
  
  % 左上角：最高气温标题
  \node[right, font=\normalsize] at (-0.3, 4.2) {最高气温 / $^\circ$C};
  % 右上角：年月
  \node[left, font=\normalsize] at (8*\xstep + 0.2, 4.8) {2026 年 1 月};
  
  % 3. 图例绘制 (与最高气温标题同一水平线 y=4.2)
  \draw[very thick, black] (3.6*\xstep, 4.2) -- (4.4*\xstep, 4.2);
  \node[right, font=\normalsize] at (4.4*\xstep, 4.2) {北京};
  \draw[very thick, dashed, black] (5.6*\xstep, 4.2) -- (6.4*\xstep, 4.2);
  \node[right, font=\normalsize] at (6.4*\xstep, 4.2) {南京};
  
  % 4. 绘制折线与数据标点
  
  % 北京数据线 (Solid Line)
  \draw[very thick, black] 
    (1*\xstep, 15*\ystep) --
    (2*\xstep, 15*\ystep) --
    (3*\xstep, 13*\ystep) --
    (4*\xstep, 10*\ystep) --
    (5*\xstep,  5*\ystep) --
    (6*\xstep,  6*\ystep) --
    (7*\xstep,  7*\ystep);
    
  % 北京标点与避障型数字标签定位
  \fill[black] (1*\xstep, 15*\ystep) circle (1.5pt) node[left, xshift=-2pt, font=\normalsize] {15};
  \fill[black] (2*\xstep, 15*\ystep) circle (1.5pt) node[below, xshift=-2pt, font=\normalsize] {15};
  \fill[black] (3*\xstep, 13*\ystep) circle (1.5pt) node[below left, font=\normalsize] {13};
  \fill[black] (4*\xstep, 10*\ystep) circle (1.5pt) node[above right, font=\normalsize] {10};
  \fill[black] (5*\xstep,  5*\ystep) circle (1.5pt) node[below left, font=\normalsize] {5};
  \fill[black] (6*\xstep,  6*\ystep) circle (1.5pt) node[below right, font=\normalsize] {6};
  \fill[black] (7*\xstep,  7*\ystep) circle (1.5pt) node[right, xshift=2pt, font=\normalsize] {7};

  % 南京数据线 (Dashed Line)
  \draw[very thick, dashed, black] 
    (1*\xstep, 25*\ystep) --
    (2*\xstep, 21*\ystep) --
    (3*\xstep, 22*\ystep) --
    (4*\xstep, 23*\ystep) --
    (5*\xstep, 25*\ystep) --
    (6*\xstep, 26*\ystep) --
    (7*\xstep, 14*\ystep);
    
  % 南京标点与避障型数字标签定位
  \fill[black] (1*\xstep, 25*\ystep) circle (1.5pt) node[above right, font=\normalsize] {25};
  \fill[black] (2*\xstep, 21*\ystep) circle (1.5pt) node[above right, font=\normalsize] {21};
  \fill[black] (3*\xstep, 22*\ystep) circle (1.5pt) node[above left, font=\normalsize] {22};
  \fill[black] (4*\xstep, 23*\ystep) circle (1.5pt) node[above left, font=\normalsize] {23};
  \fill[black] (5*\xstep, 25*\ystep) circle (1.5pt) node[above left, font=\normalsize] {25};
  \fill[black] (6*\xstep, 26*\ystep) circle (1.5pt) node[above right, font=\normalsize] {26};
  \fill[black] (7*\xstep, 14*\ystep) circle (1.5pt) node[above right, font=\normalsize] {14};

\end{tikzpicture}
\end{center}

\vspace{0.4cm}
{\small\color{blue!70!black}
\textbf{画图及逻辑控制思路解构：}\\
\textbf{1. 全域封锁的高压闭合式栅栏主轴：} 根据深度解析提取的原图坐标轴指纹，这并不是常见的三面漏风型普通直角坐标系！其本体是由 $8 \times 6$ （列向被日期及文字封锁、行向由温度刻度封死）的严密矩形封闭空间组成。为了逼真还原那一抹试卷打印特有的压抑严谨感，绘图引擎剥离了外放射线，从系统底层用 \texttt{grid} 函数铸造出了 \texttt{gray!80} 的弱灰色毛孔网纹打底，并运用厚重的纯黑色粗线条矩阵框无缝接合（即 \texttt{rectangle}）。就连最边缘的特供字幕“日期”、“最高气温”及“2026年1月”，也被我们一比一定宽死死按在了考规制定的排版锚点上！\\
\textbf{2. 极端交错路况下的全人工悬浮避障导航：} 在同一个坐标空腔里硬塞进两条缠斗交织的走势线，如何让密密麻麻的文本标签完全规避画线和相撞风险？系统为每一个悬空的孤悬分子均开启了独立的坐标偏置锚泊模块（如 \texttt{above left/right} 定向锁定区）。以实线狂跌区的北京队 $4$ 日节点为例（数字 $10$），为了给左翼的极速下坠陡坡让出缓冲视觉带，系统将其极其刁钻地甩挂到了右偏上方（\texttt{above right}）的安全防空洞中；而反观虚线区的南京队大牛市主升浪（$22-25$ 这整整三天的震荡抬升），系统引擎又强行将数字牌阵列集体向左方冷空域进行了倒板退让（\texttt{above left}），生怕其刮擦到粗黑上扬的骨架主体。没有任何数值重叠相撞，在极端密集的重装战区展现出极致的图纸解剖艺术感！
}

%% ==============================
%% 第47题（补充题）：复式折线统计图——水费电费
%% ==============================
\newpage



%% ==============================
%% 第47题（补充题）：复式折线统计图——水费电费
%% ==============================
\newpage

\section{解答：复式折线统计图——水电费支出}

\vspace{0.6cm}

\textbf{\LARGE 38.} \textbf{收集你家去年下半年水、电费支出数据（精确到整元数），并完成下面的折线统计图。}

\vspace{0.4cm}
\textbf{解：}

\textbf{（1）题目分析：}

题目要求收集去年下半年（7 月至 12 月）的水费和电费数据，绘制复式折线统计图。根据参考答案，数据如下：

\vspace{0.2cm}
\begin{center}
\renewcommand{\arraystretch}{1.5}
\begin{tabular}{c|c|c|c|c|c|c}
\toprule[1.2pt]
月份 & 七月 & 八月 & 九月 & 十月 & 十一月 & 十二月 \\
\hline
水费 / 元 & 100 & 150 & 50 & 60 & 80 & 120 \\
\hline
电费 / 元 & 450 & 500 & 350 & 300 & 350 & 400 \\
\bottomrule[1.2pt]
\end{tabular}
\end{center}

\vspace{0.4cm}
\textbf{（2）折线统计图绘制：}

\begin{center}
{\textbf{\large \underline{小明}家去年下半年水、电费支出情况统计图}}
\end{center}

\vspace{0.2cm}
\begin{center}
\begin{tikzpicture}[>=Stealth, line join=round]
  % 核心缩放参数
  \def\xstep{1.5}  % X轴每月间距
  \def\ystep{0.01} % Y轴每元对应长度，500元对应5cm

  % 1. 绘制网格与外框
  % X轴共7格(0~7)，Y轴共10格(0~500，每50一格)
  \draw[gray!60, thin, xscale=\xstep, yscale=50*\ystep] (0,0) grid (7,10);
  \draw[thick, black] (0,0) rectangle (7*\xstep, 500*\ystep);

  % 2. Y轴刻度标签 (每100标一个)
  \foreach \y in {0, 50, 100, 150, 200, 250, 300, 350, 400, 450, 500} {
      \node[left, font=\small] at (0, \y*\ystep) {\y};
  }

  % 3. X轴月份标签
  \foreach \x/\label in {1/七, 2/八, 3/九, 4/十, 5/十一, 6/十二} {
      \node[below, font=\normalsize] at (\x*\xstep, -0.15) {\label};
  }
  \node[below, font=\normalsize] at (7*\xstep, -0.15) {月份};

  % 4. 图表顶部信息 (绝对Y坐标，500*0.01=5.0cm为顶部)
  % 图例行 y=5.6，年份行 y=6.2
  \node[right, font=\normalsize] at (-0.3, 5.6) {数量/元};
  \node[left, font=\normalsize] at (7*\xstep + 0.2, 6.2) {2026 年 1 月};

  % 图例
  \draw[very thick, black] (3.0*\xstep, 5.6) -- (3.8*\xstep, 5.6);
  \node[right, font=\normalsize] at (3.8*\xstep, 5.6) {水费};
  \draw[very thick, dashed, black] (5.0*\xstep, 5.6) -- (5.8*\xstep, 5.6);
  \node[right, font=\normalsize] at (5.8*\xstep, 5.6) {电费};

  % 5. 水费折线 (实线)
  \draw[very thick, black]
    (1*\xstep, 100*\ystep) --
    (2*\xstep, 150*\ystep) --
    (3*\xstep,  50*\ystep) --
    (4*\xstep,  60*\ystep) --
    (5*\xstep,  80*\ystep) --
    (6*\xstep, 120*\ystep);

  % 水费标点与数值标注
  \fill[black] (1*\xstep, 100*\ystep) circle (1.5pt) node[below, font=\small, text=red] {100};
  \fill[black] (2*\xstep, 150*\ystep) circle (1.5pt) node[above right, font=\small, text=red] {150};
  \fill[black] (3*\xstep,  50*\ystep) circle (1.5pt) node[below, font=\small, text=red] {50};
  \fill[black] (4*\xstep,  60*\ystep) circle (1.5pt) node[below, font=\small, text=red] {60};
  \fill[black] (5*\xstep,  80*\ystep) circle (1.5pt) node[below, font=\small, text=red] {80};
  \fill[black] (6*\xstep, 120*\ystep) circle (1.5pt) node[below right, font=\small, text=red] {120};

  % 6. 电费折线 (虚线)
  \draw[very thick, dashed, black]
    (1*\xstep, 450*\ystep) --
    (2*\xstep, 500*\ystep) --
    (3*\xstep, 350*\ystep) --
    (4*\xstep, 300*\ystep) --
    (5*\xstep, 350*\ystep) --
    (6*\xstep, 400*\ystep);

  % 电费标点与数值标注
  \fill[black] (1*\xstep, 450*\ystep) circle (1.5pt) node[above left, font=\small, text=red] {450};
  \fill[black] (2*\xstep, 500*\ystep) circle (1.5pt) node[above right, font=\small, text=red] {500};
  \fill[black] (3*\xstep, 350*\ystep) circle (1.5pt) node[above right, font=\small, text=red] {350};
  \fill[black] (4*\xstep, 300*\ystep) circle (1.5pt) node[above, font=\small, text=red] {300};
  \fill[black] (5*\xstep, 350*\ystep) circle (1.5pt) node[above left, font=\small, text=red] {350};
  \fill[black] (6*\xstep, 400*\ystep) circle (1.5pt) node[above right, font=\small, text=red] {400};

\end{tikzpicture}
\end{center}

\vspace{0.3cm}
\textbf{（3）从图中获取的信息：}

\begin{itemize}
    \item 电费支出远高于水费支出，每月电费约为水费的 $3 \sim 6$ 倍。
    \item 八月电费最高（$500$ 元），九月水费最低（$50$ 元）。
    \item 夏季（七、八月）水电费均较高，秋季（九、十月）有所回落。
\end{itemize}

\vspace{0.4cm}
{\small\color{blue!70!black}
\textbf{绘图原理与步骤：}\\
\textbf{1. 坐标系结构：} 采用 $7 \times 10$ 封闭矩形网格。X 轴 7 格对应 6 个月份及末尾"月份"标签，Y 轴 10 格对应 $0$--$500$ 元（每格 $50$ 元）。缩放参数 \texttt{\textbackslash{}xstep=1.5}（每月间距 1.5cm），\texttt{\textbackslash{}ystep=0.01}（每元 0.01cm，即 $500$ 元 $= 5$cm）。\\
\textbf{2. 图例与标题定位：} 使用绝对 Y 坐标。图表顶部为 $500 \times 0.01 = 5.0$cm，图例行固定于 $y=5.6$cm，年份行固定于 $y=6.2$cm，确保不与网格重叠。\\
\textbf{3. 双线绘制：} 水费使用实线（\texttt{very thick, black}），电费使用虚线（\texttt{very thick, dashed, black}）。两组数据分离度大（水费 $50$--$150$，电费 $300$--$500$），无交叉冲突。\\
\textbf{4. 标注策略：} 数值标签使用红色（\texttt{text=red}）突出显示。水费数据较低，标签多置于下方（\texttt{below}）；电费数据较高，标签多置于上方（\texttt{above left/right}），根据折线走势方向选择左右偏移以避免遮挡连线。
}

%% ==============================
%% 第48题（补充题）：10天最高最低气温统计表与折线图
%% ==============================
\newpage



%% ==============================
%% 第48题（补充题）：10天最高最低气温统计表与折线图
%% ==============================
\newpage

\section{解答：10天气温调查统计表与复式折线统计图}

\vspace{0.6cm}

\textbf{\LARGE 39.} \textbf{调查了解本地近 10 天的最高气温和最低气温，并分别完成下面的统计表和折线统计图。}

\vspace{0.4cm}
\textbf{解：}

\textbf{（1）统计表：}

\begin{center}
{\textbf{\large \underline{上海}地区\underline{\,1\,}月\underline{\,1\,}日至\underline{\,1\,}月\underline{\,10\,}日气温情况统计表}}
\end{center}

%\vspace{-0.4cm}
\begin{center}
\renewcommand{\arraystretch}{1.4}
\begin{tabular}{c|c|c|c|c|c|c|c|c|c|c}
\multicolumn{11}{r}{2026 年 1 月} \\[2pt]
\noalign{\hrule height 1.2pt}
日\quad 期 & 1.1 & 1.2 & 1.3 & 1.4 & 1.5 & 1.6 & 1.7 & 1.8 & 1.9 & 1.10 \\
\hline
最高气温/$^\circ$C & 14 & 14 & 11 & 8 & 12 & 14 & 17 & 17 & 16 & 15 \\
\hline
最低气温/$^\circ$C & 6 & 9 & 2 & 1 & 1 & 1 & 4 & 6 & 8 & 7 \\
\noalign{\hrule height 1.2pt}
\end{tabular}
\end{center}

\vspace{0.6cm}
\textbf{（2）折线统计图：}

\begin{center}
{\textbf{\large \underline{上海}地区\underline{\,1\,}月\underline{\,1\,}日至\underline{\,1\,}月\underline{\,10\,}日气温情况统计图}}
\end{center}

\vspace{0.2cm}
\begin{center}
\begin{tikzpicture}[>=Stealth, line join=round]
  % 缩放参数
  \def\xstep{1.0}  % X轴每天间距 1.0cm
  \def\ystep{0.14}  % Y轴每度 0.14cm => 35度=4.9cm

  % 1. 网格与外框
  % X轴11格(0~11)，Y轴7格(0~35，每5一格)
  \draw[gray!60, thin, xscale=\xstep, yscale=5*\ystep] (0,0) grid (11,7);
  \draw[thick, black] (0,0) rectangle (11*\xstep, 35*\ystep);

  % 2. Y轴标签
  \foreach \y in {0, 5, 10, 15, 20, 25, 30, 35} {
      \node[left, font=\small] at (0, \y*\ystep) {\y};
  }

  % 3. X轴标签
  \foreach \x in {1,...,10} {
      \node[below, font=\small] at (\x*\xstep, -0.15) {\x};
  }
  \node[below, font=\normalsize] at (11*\xstep, -0.15) {第几天};

  % 4. 顶部信息 (绝对Y坐标: 35*0.14=4.9cm为顶部)
  % 图例行 y=5.5，年份行 y=6.0
  \node[right, font=\normalsize] at (-0.3, 5.5) {气温/$^\circ$C};
  \node[left, font=\normalsize] at (11*\xstep + 0.2, 6.0) {2026 年 1 月};

  % 图例
  \draw[very thick, black] (3.5*\xstep, 5.5) -- (4.5*\xstep, 5.5);
  \node[right, font=\normalsize] at (4.5*\xstep, 5.5) {最高气温};
  \draw[very thick, dashed, black] (6.5*\xstep, 5.5) -- (7.5*\xstep, 5.5);
  \node[right, font=\normalsize] at (7.5*\xstep, 5.5) {最低气温};

  % 5. 最高气温折线 (实线)
  % 数据: 14, 14, 11, 8, 12, 14, 17, 17, 16, 15
  \draw[very thick, black]
    ( 1*\xstep, 14*\ystep) --
    ( 2*\xstep, 14*\ystep) --
    ( 3*\xstep, 11*\ystep) --
    ( 4*\xstep,  8*\ystep) --
    ( 5*\xstep, 12*\ystep) --
    ( 6*\xstep, 14*\ystep) --
    ( 7*\xstep, 17*\ystep) --
    ( 8*\xstep, 17*\ystep) --
    ( 9*\xstep, 16*\ystep) --
    (10*\xstep, 15*\ystep);

  % 最高气温标注
  \fill[black] ( 1*\xstep, 14*\ystep) circle (1.5pt) node[above left, font=\small, text=red] {14};
  \fill[black] ( 2*\xstep, 14*\ystep) circle (1.5pt) node[above, font=\small, text=red] {14};
  \fill[black] ( 3*\xstep, 11*\ystep) circle (1.5pt) node[above, font=\small, text=red] {11};
  \fill[black] ( 4*\xstep,  8*\ystep) circle (1.5pt) node[above right, font=\small, text=red] {8};
  \fill[black] ( 5*\xstep, 12*\ystep) circle (1.5pt) node[above right, font=\small, text=red] {12};
  \fill[black] ( 6*\xstep, 14*\ystep) circle (1.5pt) node[above, font=\small, text=red] {14};
  \fill[black] ( 7*\xstep, 17*\ystep) circle (1.5pt) node[above left, font=\small, text=red] {17};
  \fill[black] ( 8*\xstep, 17*\ystep) circle (1.5pt) node[above right, font=\small, text=red] {17};
  \fill[black] ( 9*\xstep, 16*\ystep) circle (1.5pt) node[above, font=\small, text=red] {16};
  \fill[black] (10*\xstep, 15*\ystep) circle (1.5pt) node[above right, font=\small, text=red] {15};

  % 6. 最低气温折线 (虚线)
  % 数据: 6, 9, 2, 1, 1, 1, 4, 6, 8, 7
  \draw[very thick, dashed, black]
    ( 1*\xstep,  6*\ystep) --
    ( 2*\xstep,  9*\ystep) --
    ( 3*\xstep,  2*\ystep) --
    ( 4*\xstep,  1*\ystep) --
    ( 5*\xstep,  1*\ystep) --
    ( 6*\xstep,  1*\ystep) --
    ( 7*\xstep,  4*\ystep) --
    ( 8*\xstep,  6*\ystep) --
    ( 9*\xstep,  8*\ystep) --
    (10*\xstep,  7*\ystep);

  % 最低气温标注
  \fill[black] ( 1*\xstep,  6*\ystep) circle (1.5pt) node[below left, font=\small, text=red] {6};
  \fill[black] ( 2*\xstep,  9*\ystep) circle (1.5pt) node[below, font=\small, text=red] {9};
  \fill[black] ( 3*\xstep,  2*\ystep) circle (1.5pt) node[below, font=\small, text=red] {2};
  \fill[black] ( 4*\xstep,  1*\ystep) circle (1.5pt) node[below right, font=\small, text=red] {1};
  \fill[black] ( 5*\xstep,  1*\ystep) circle (1.5pt) node[below, font=\small, text=red] {1};
  \fill[black] ( 6*\xstep,  1*\ystep) circle (1.5pt) node[below left, font=\small, text=red] {1};
  \fill[black] ( 7*\xstep,  4*\ystep) circle (1.5pt) node[below right, font=\small, text=red] {4};
  \fill[black] ( 8*\xstep,  6*\ystep) circle (1.5pt) node[below right, font=\small, text=red] {6};
  \fill[black] ( 9*\xstep,  8*\ystep) circle (1.5pt) node[below, font=\small, text=red] {8};
  \fill[black] (10*\xstep,  7*\ystep) circle (1.5pt) node[below right, font=\small, text=red] {7};

\end{tikzpicture}
\end{center}

\vspace{0.4cm}
{\small\color{blue!70!black}
\textbf{绘图原理与步骤：}\\
\textbf{1. 坐标系结构：} 采用 $11 \times 7$ 封闭矩形网格。X 轴 11 格对应 10 天及末尾"第几天"标签，Y 轴 7 格对应 $0$--$35$$^\circ$C（每格 $5$$^\circ$C）。缩放参数 \texttt{\textbackslash{}xstep=1.0}，\texttt{\textbackslash{}ystep=0.14}（每度 0.14cm，$35$ 度 $= 4.9$cm）。\\
\textbf{2. 图例与标题定位：} 绝对 Y 坐标。图表顶部 $35 \times 0.14 = 4.9$cm，图例行 $y=5.5$cm，年份行 $y=6.0$cm。\\
\textbf{3. 双线绘制：} 最高气温实线，最低气温虚线。两组数据无交叉（最高 $8$--$17$，最低 $1$--$9$），标签分设上下方避碰。\\
\textbf{4. 标注策略：} 红色文字标注。最高气温标签置于上方，最低气温标签置于下方，根据线段走向选择 left/right 偏移防撞。
}

%% ==============================
%% 第49题（补充题）：条形统计图补全与扇形统计图
%% ==============================
\newpage



%% ==============================
%% 第49题（补充题）：条形统计图补全与扇形统计图
%% ==============================
\newpage

\section{解答：条形统计图与扇形统计图综合}

\vspace{0.6cm}

\textbf{\LARGE 40.} \textbf{（8 分）某校开展了"文明校园"活动周，活动周设置了"A：文明礼仪，B：生态环境，C：交通安全，D：卫生保洁"四个主题活动，每个学生限选一个主题参与。为了了解活动开展情况，学校随机抽取了部分学生进行调查，并根据调查结果绘制了如图所示的不完整的条形统计图和扇形统计图。}

\vspace{0.4cm}
\textbf{解：}

\textbf{（1）解题分析：}

由扇形统计图知 A 占 $25\%$，条形图中 A $= 15$ 人，故总人数 $= 15 \div 25\% = 60$ 人。

C 主题人数 $= 60 - 15 - 18 - 9 = 18$ 人。

\vspace{0.2cm}
\textbf{（2）补全条形统计图：}

\begin{center}
\begin{tikzpicture}[>=Stealth]
  % 缩放参数
  \def\xstep{1.2}   % 每组柱子间距
  \def\ystep{0.22}   % Y轴每人的高度，18人=3.96cm
  \def\barw{0.7}     % 条形柱宽度

  % 1. Y轴与X轴
  \draw[thick, black, ->] (0,0) -- (0, 20*\ystep + 0.5) node[left, font=\normalsize] {人数};
  \draw[thick, black, ->] (0,0) -- (5*\xstep + 0.3, 0);

  % 2. Y轴刻度标签（仅标签，不画虚线）
  \foreach \y in {3, 6, 9, 12, 15, 18} {
      \node[left, font=\small] at (0, \y*\ystep) {\y};
  }
  % 0刻度
  \node[left, font=\small] at (0, 0) {0};

  % 3. 绘制条形柱（带浅灰填充）
  % A = 15
  \fill[gray!25] (1*\xstep - \barw/2, 0) rectangle (1*\xstep + \barw/2, 15*\ystep);
  \draw[thick, black] (1*\xstep - \barw/2, 0) rectangle (1*\xstep + \barw/2, 15*\ystep);
  \node[above, font=\small\bfseries] at (1*\xstep, 15*\ystep) {15};

  % B = 18
  \fill[gray!25] (2*\xstep - \barw/2, 0) rectangle (2*\xstep + \barw/2, 18*\ystep);
  \draw[thick, black] (2*\xstep - \barw/2, 0) rectangle (2*\xstep + \barw/2, 18*\ystep);
  \node[above, font=\small\bfseries] at (2*\xstep, 18*\ystep) {18};

  % C = 18（补全项，红色标注）
  \fill[gray!25] (3*\xstep - \barw/2, 0) rectangle (3*\xstep + \barw/2, 18*\ystep);
  \draw[thick, black] (3*\xstep - \barw/2, 0) rectangle (3*\xstep + \barw/2, 18*\ystep);
  \node[above, font=\small\bfseries, text=red] at (3*\xstep, 18*\ystep) {18};

  % D = 9
  \fill[gray!25] (4*\xstep - \barw/2, 0) rectangle (4*\xstep + \barw/2, 9*\ystep);
  \draw[thick, black] (4*\xstep - \barw/2, 0) rectangle (4*\xstep + \barw/2, 9*\ystep);
  \node[above, font=\small\bfseries] at (4*\xstep, 9*\ystep) {9};

  % 4. 水平虚线——从Y轴延伸到对应条形柱左边缘即截止，在柱上层
  % y=9 → D柱 (x=4*xstep)
  \draw[dashed, gray!70] (0, 9*\ystep) -- (4*\xstep - \barw/2, 9*\ystep);
  % y=15 → A柱 (x=1*xstep)
  \draw[dashed, gray!70] (0, 15*\ystep) -- (1*\xstep - \barw/2, 15*\ystep);
  % y=18 → B柱 (x=2*xstep)
  \draw[dashed, gray!70] (0, 18*\ystep) -- (2*\xstep - \barw/2, 18*\ystep);

  % 5. X轴标签
  \foreach \x/\label in {1/A, 2/B, 3/C, 4/D} {
      \node[below, font=\normalsize] at (\x*\xstep, -0.15) {\label};
  }
  \node[below, font=\normalsize] at (5*\xstep, -0.15) {主题};

\end{tikzpicture}
\end{center}

\vspace{0.4cm}
\textbf{（3）各问作答：}

\begin{itemize}
    \item[\textbf{(1)}] 本次随机调查的学生人数是 $\underline{\textcolor{red}{\,60\,}}$ 人。

    （$15 \div 25\% = 60$）

    \item[\textbf{(2)}] 条形统计图已补全（C 主题 $= 60 - 15 - 18 - 9 = 18$ 人）。

    \item[\textbf{(3)}] B 主题对应扇形的圆心角 $= 360^\circ \times \dfrac{18}{60} = \underline{\textcolor{red}{\,108\,}}^\circ$。

    \item[\textbf{(4)}] 估计参与"生态环境"主题的学生人数 $= 1800 \times \dfrac{18}{60} = \underline{\textcolor{red}{\,540\,}}$ 人。
\end{itemize}

\vspace{0.4cm}
{\small\color{blue!70!black}
\textbf{绘图原理与步骤：}\\
\textbf{1. 坐标系：} 开放式直角坐标系（Y 轴带箭头）。Y 轴 $0$--$18$，每 $3$ 人一条水平虚线（\texttt{dashed, gray!70}）。X 轴放置 A/B/C/D 及"主题"标签。缩放参数 \texttt{\textbackslash{}xstep=1.2}，\texttt{\textbackslash{}ystep=0.22}，\texttt{\textbackslash{}barw=0.7}。\\
\textbf{2. 条形柱：} \texttt{\textbackslash{}fill[gray!25]} 填充浅灰，\texttt{\textbackslash{}draw[thick]} 黑色描边。补全项 C 的数值标签用红色（\texttt{text=red}）。\\
\textbf{3. 数据推导：} A 占 $25\% = 15$ 人 $\Rightarrow$ 总数 $60$ 人。C $= 60 - 15 - 18 - 9 = 18$。B 圆心角 $= 360^\circ \times 30\% = 108^\circ$。
}

%% ==============================
%% 第50题（补充题）：代数拼图——矩形纸片拼接
%% ==============================
\newpage



%% ==============================
%% 第51题（补充题）：图书类别条形统计图补全
%% ==============================
\newpage

\section{解答：图书类别条形统计图补全与扇形统计图}

\vspace{0.6cm}

\textbf{\LARGE 42.} \textbf{（8 分）某中学准备购进一批图书供学生阅读，为了合理配备各类图书，从全体学生中随机抽取了部分学生进行了问卷调查。问卷设置了五种选项：A"艺术类"，B"文学类"，C"科普类"，D"体育类"，E"其他类"。每名学生必须且只能选择其中最喜爱的一类图书，将调查结果整理绘制成如下两幅不完整的统计图。}

\vspace{0.4cm}
\textbf{解：}

\textbf{（1）解题分析：}

由扇形统计图知 B 占 $20\%$，条形图中 B $= 20$ 人，故抽样总数 $= 20 \div 20\% = 100$ 人。

D 的人数 $= 100 - 10 - 20 - 40 - 5 = 25$ 人。

\vspace{0.2cm}
\textbf{（2）补全条形统计图：}

\begin{center}
{\textbf{学生最喜爱图书类别的人数条形统计图}}
\end{center}

\vspace{0.1cm}
\begin{center}
\begin{tikzpicture}[>=Stealth]
  % 缩放参数
  \def\xstep{1.2}   % 每组间距
  \def\ystep{0.12}   % Y轴每人高度，40人=4.8cm
  \def\barw{0.7}     % 柱宽

  % 1. Y轴与X轴
  \draw[thick, black, ->] (0,0) -- (0, 42*\ystep + 0.5) node[left, font=\normalsize] {人数/名};
  \draw[thick, black, ->] (0,0) -- (6*\xstep + 0.3, 0);

  % 2. Y轴刻度标签（仅标签，不画虚线）
  \foreach \y in {5, 10, 15, 20, 25, 30, 35, 40} {
      \node[left, font=\small] at (0, \y*\ystep) {\y};
  }
  \node[left, font=\small] at (0, 0) {0};

  % 3. 条形柱（先绘制，虚线后绘制以覆盖在上层）
  % A = 10（艺术类）
  \fill[gray!35] (1*\xstep - \barw/2, 0) rectangle (1*\xstep + \barw/2, 10*\ystep);
  \draw[thick, black] (1*\xstep - \barw/2, 0) rectangle (1*\xstep + \barw/2, 10*\ystep);
  \node[above, font=\small\bfseries] at (1*\xstep, 10*\ystep) {10};

  % B = 20（文学类）
  \fill[gray!35] (2*\xstep - \barw/2, 0) rectangle (2*\xstep + \barw/2, 20*\ystep);
  \draw[thick, black] (2*\xstep - \barw/2, 0) rectangle (2*\xstep + \barw/2, 20*\ystep);
  \node[above, font=\small\bfseries] at (2*\xstep, 20*\ystep) {20};

  % C = 40（科普类）
  \fill[gray!35] (3*\xstep - \barw/2, 0) rectangle (3*\xstep + \barw/2, 40*\ystep);
  \draw[thick, black] (3*\xstep - \barw/2, 0) rectangle (3*\xstep + \barw/2, 40*\ystep);
  \node[above, font=\small\bfseries] at (3*\xstep, 40*\ystep) {40};

  % D = 25（体育类，补全项，红色标注）
  \fill[gray!35] (4*\xstep - \barw/2, 0) rectangle (4*\xstep + \barw/2, 25*\ystep);
  \draw[thick, black] (4*\xstep - \barw/2, 0) rectangle (4*\xstep + \barw/2, 25*\ystep);
  \node[above, font=\small\bfseries, text=red] at (4*\xstep, 25*\ystep) {25};

  % E = 5（其他类）
  \fill[gray!35] (5*\xstep - \barw/2, 0) rectangle (5*\xstep + \barw/2, 5*\ystep);
  \draw[thick, black] (5*\xstep - \barw/2, 0) rectangle (5*\xstep + \barw/2, 5*\ystep);
  \node[above, font=\small\bfseries] at (5*\xstep, 5*\ystep) {5};

  % 4. 水平虚线——从Y轴延伸到对应条形柱左边缘即截止
  % y=10 → A柱 (x=1*xstep)
  \draw[dashed, gray!70] (0, 10*\ystep) -- (1*\xstep - \barw/2, 10*\ystep);
  % y=20 → B柱 (x=2*xstep)
  \draw[dashed, gray!70] (0, 20*\ystep) -- (2*\xstep - \barw/2, 20*\ystep);
  % y=40 → C柱 (x=3*xstep)
  \draw[dashed, gray!70] (0, 40*\ystep) -- (3*\xstep - \barw/2, 40*\ystep);
  % y=25 → D柱 (x=4*xstep)
  \draw[dashed, gray!70] (0, 25*\ystep) -- (4*\xstep - \barw/2, 25*\ystep);
  % y=5 → E柱 (x=5*xstep)
  \draw[dashed, gray!70] (0, 5*\ystep) -- (5*\xstep - \barw/2, 5*\ystep);

  % 5. X轴标签
  \foreach \x/\label in {1/A, 2/B, 3/C, 4/D, 5/E} {
      \node[below, font=\normalsize] at (\x*\xstep, -0.15) {\label};
  }
  \node[below, font=\normalsize] at (6*\xstep, -0.15) {类别};

\end{tikzpicture}
\end{center}

\vspace{0.4cm}
\textbf{（3）各问作答：}

\begin{itemize}
    \item[\textbf{(1)}] 抽样总数 $= 20 \div 20\% = \underline{\textcolor{red}{\,100\,}}$ 人。

    D 的人数 $= 100 - 10 - 20 - 40 - 5 = 25$ 人，条形统计图已补全。

    \item[\textbf{(2)}] A"艺术类"所对应的圆心角 $= 360^\circ \times \dfrac{10}{100} = \underline{\textcolor{red}{\,36\,}}^\circ$。

    \item[\textbf{(3)}] 估计喜爱 C"科普类"图书的学生人数 $= 1200 \times \dfrac{40}{100} = \underline{\textcolor{red}{\,480\,}}$ 人。
\end{itemize}

\vspace{0.4cm}
{\small\color{blue!70!black}
\textbf{绘图原理与步骤：}\\
\textbf{1. 坐标系：} 开放式直角坐标系（双轴带箭头）。Y 轴 $0$--$40$，每 $5$ 人一条水平虚线。X 轴 5 组类别 + "类别"标签。缩放参数 \texttt{\textbackslash{}xstep=1.2}，\texttt{\textbackslash{}ystep=0.12}，\texttt{\textbackslash{}barw=0.7}。\\
\textbf{2. 条形柱：} \texttt{gray!35} 浅灰填充 + 黑色描边。补全项 D 的数值用红色标注。\\
\textbf{3. 数据推导：} B 占 $20\% = 20$ 人 $\Rightarrow$ 总数 $100$。D $= 100 - 10 - 20 - 40 - 5 = 25$。A 圆心角 $= 360^\circ \times 10\% = 36^\circ$。
}

%% ==============================
%% 第52题（补充题）：Rt△ABC 中线、尺规作图与菱形证明
%% ==============================
\newpage



%% ==============================
%% 第77题（补充题）：统计表与统计图（全套数据链）
%% ==============================
\newpage

\section{解答：数据的收集与整理}

\vspace{0.6cm}

\textbf{\LARGE 68.} \textbf{（第 3 题）李梅调查了本班同学最喜爱的一项体育活动，调查结果如下。}\\

% --- 计票原始表格 ---
% 使用已有的 \zhengII/\zhengIII/\zhengIV 宏绘制不完整正字笔画

\begin{table}[H]
\centering
\renewcommand{\arraystretch}{1.5}
\begin{tabular}{|c|c|c|c|c|}
\hline
\makebox[1.8cm]{跳\quad 绳} & \makebox[1.8cm]{滑\quad 冰} & \makebox[1.8cm]{踢足球} & \makebox[1.8cm]{打乒乓球} & \makebox[1.8cm]{游\quad 泳} \\
\hline
正 & 正正\,\zhengII & 正正 & \zhengIV & 正\,\zhengIII \\
\hline
\end{tabular}
\end{table}

\vspace{0.4cm}
\textbf{根据调查结果完成下面的统计表和统计图。}

\vspace{0.6cm}
\textbf{解：}

由计票结果可知各项目人数为：
跳绳 $5$ 人；滑冰 $12$ 人；踢足球 $10$ 人；打乒乓球 $4$ 人；游泳 $8$ 人。\\
总人数为 $5 + 12 + 10 + 4 + 8 = 39$ 人。

\vspace{0.4cm}

% --- 答案统计表 ---
\begin{center}
\textbf{\large 某班同学最喜爱的一项体育活动统计表}

\vspace{0.2cm}
\hfill \makebox[2.5cm][r]{{\color{blue}3} \ 月 \ {\color{blue}9} \ 日}
\end{center}

\vspace{-0.4cm}
\begin{table}[H]
\centering
\renewcommand{\arraystretch}{1.5}
% 表格样式设定：左右两边无竖线框，内含列分割，上下有横线
\begin{tabular}{c|c|c|c|c|c|c}
\hline
项\quad 目 & 合\quad 计 & 跳\quad 绳 & 滑\quad 冰 & 踢足球 & 打乒乓球 & 游\quad 泳 \\
\hline
人\quad 数 & \textcolor{blue}{39} & \textcolor{blue}{5} & \textcolor{blue}{12} & \textcolor{blue}{10} & \textcolor{blue}{4} & \textcolor{blue}{8} \\
\hline
\end{tabular}
\end{table}

\vspace{0.8cm}

% --- 答案统计图（条形） ---
\begin{center}
\textbf{\large 某班同学最喜爱的一项体育活动统计图}

\vspace{0.5cm}
\begin{tikzpicture}[x=0.8cm, y=0.35cm]
  % 标题附属文字
  \node[above] at (0, 14.5) {数量/人};
  \node[above] at (11, 14.5) {{\color{blue}3} \ 月 \ {\color{blue}9} \ 日};

  % 绘制主网格（共11列，7行大格）
  \draw[black, line width=0.8pt] (0,0) rectangle (11,14);
  
  % 内部横线（每行代表2人，真实高度标尺为14，所以画到12即可）
  \foreach \y in {2,4,6,8,10,12} {
    \draw[black, line width=0.4pt] (0,\y) -- (11,\y);
  }
  
  % 内部竖线（共10条内部竖线，划分出11个空格）
  \foreach \x in {1,2,...,10} {
    \draw[black, line width=0.4pt] (\x,0) -- (\x,14);
  }
  
  % 左侧 Y 轴刻度文字
  \foreach \y in {0,2,4,6,8,10,12,14} {
    \node[left, font=\small] at (0,\y) {\y};
  }
  
  % 底部 X 轴项目类别名称（位于柱子中间）
  \node[below, font=\small] at (1.5, -0.3) {跳绳};
  \node[below, font=\small] at (3.5, -0.3) {滑冰};
  \node[below, font=\small] at (5.5, -0.3) {踢足球};
  \node[below, font=\small] at (7.5, -0.3) {打乒乓球};
  \node[below, font=\small] at (9.5, -0.3) {游泳};
  
  \node[below, right, font=\small] at (11, -0.3) {项目};

  % 绘制数据蓝柱子（完美对齐王哥/网格背景边界）
  % 设间隔为空白格。第1个矩形位于(1,0)-(2,y)，第2个位于(3,0)-(4,y)...
  \fill[cyan!70!blue] (1, 0) rectangle (2, 5);
  \node[above, font=\large] at (1.5, 5) {5};
  
  \fill[cyan!70!blue] (3, 0) rectangle (4, 12);
  \node[above, font=\large] at (3.5, 12) {12};
  
  \fill[cyan!70!blue] (5, 0) rectangle (6, 10);
  \node[above, font=\large] at (5.5, 10) {10};
  
  \fill[cyan!70!blue] (7, 0) rectangle (8, 4);
  \node[above, font=\large] at (7.5, 4) {4};
  
  \fill[cyan!70!blue] (9, 0) rectangle (10, 8);
  \node[above, font=\large] at (9.5, 8) {8};

\end{tikzpicture}
\end{center}

\vspace{0.4cm}
{\small\color{blue!70!black}
\textbf{绘图原理与步骤：}\\
\textbf{1. 计票宏复用与字体拟合：} 沿用文档头部的 \texttt{\textbackslash zhengII} 至 \texttt{\textbackslash zhengIV} 宏，精准嵌入计票表格残缺汉字笔画。\\
\textbf{2. 方阵化条形图格栅排布：} 解析原题照片可以提取出一个极其隐蔽的排版内核：它根本不是 $7$ 列均匀分布柱，而是一张标准的“隔一画一” $11$ 列方形网格（包含 $11$ 个方格列）。代码依据此规则，使用 \texttt{\textbackslash foreach \textbackslash x in \{1,2,...,10\}} 布下硬核钢线网格；所有的蓝色直方图块（如跳绳在 \texttt{x=1} 到 \texttt{x=2} 之间）均使用绝对贴边填充，彻底抹除空隙 \texttt{margin}，严格落实让柱状图“强制对齐网格（即王哥）背景”的视觉标准。
}

%% ==============================
%% 第78题（补充题）：横向条形图与数据转换
%% ==============================
\newpage



%% ==============================
%% 第78题（补充题）：横向条形图与数据转换
%% ==============================
\newpage

\section{解答：横向统计图的识别与制备}

\vspace{0.6cm}

\textbf{\LARGE 69.} \textbf{（第 1 题）东阳小学合唱队女生的身高情况记录如下。（单位：厘米）}

\vspace{0.4cm}
% --- 原始数据表格 ---
\begin{table}[H]
\centering
\renewcommand{\arraystretch}{1.5}
\begin{tabular}{|c|c|c|c|c|c|c|c|}
\hline
姓\quad 名 & 身\quad 高 & 姓\quad 名 & 身\quad 高 & 姓\quad 名 & 身\quad 高 & 姓\quad 名 & 身\quad 高 \\
\hline
吴\quad 凡 & $143$ & 王思雨 & $138$ & 陈一洋 & $132$ & 吴\quad 菁 & $138$ \\
\hline
刘一辰 & $139$ & 刘\quad 婧 & $147$ & 刘\quad 丽 & $140$ & 陈君好 & $146$ \\
\hline
王\quad 晓 & $145$ & 尹\quad 鸣 & $139$ & 朱平凡 & $139$ & 杨\quad 柳 & $143$ \\
\hline
葛\quad 楠 & $135$ & 卢小菲 & $141$ & 吴\quad 婷 & $148$ & 夏晴天 & $136$ \\
\hline
凌雨桐 & $144$ & 沈语嘉 & $144$ & 赵\quad 婕 & $150$ & 张\quad 焱 & $142$ \\
\hline
\end{tabular}
\end{table}

\textbf{身高低于 140 厘米的适合穿小号服装，身高 140$\sim$145 厘米的适合穿中号服装，身高超过 145 厘米的适合穿大号服装。根据上表完成下面的统计图。}\\
\textbf{根据完成的统计图填空：小号服装要准备(\qquad)套，中号服装比大号服装要多准备(\qquad)套。}

\vspace{0.8cm}
\textbf{解：}

由上表数据逐一筛选统计：
\begin{itemize}
    \item \textbf{低于 140 厘米（小号）}：王思雨(138)、陈一洋(132)、吴菁(138)、刘一辰(139)、尹鸣(139)、朱平凡(139)、葛楠(135)、夏晴天(136)，共计 $8$ 人。
    \item \textbf{140$\sim$145 厘米（中号）}：吴凡(143)、刘丽(140)、王晓(145)、杨柳(143)、卢小菲(141)、凌雨桐(144)、沈语嘉(144)、张焱(142)，共计 $8$ 人。
    \item \textbf{超过 145 厘米（大号）}：刘婧(147)、陈君好(146)、吴婷(148)、赵婕(150)，共计 $4$ 人。
\end{itemize}

计算得出：小号服装要 $8$ 套。中号服装 $8$ 套，大号服装 $4$ 套，中号比大号多 $8 - 4 = 4$ 套。

%\vspace{0.4cm}
\newpage
% --- 横向条形统计图 ---
\begin{center}
\textbf{东阳小学合唱队女生身高情况统计图}

\vspace{0.6cm}
\begin{tikzpicture}[x=1.8cm, y=0.71cm]
  % 标题附属文字
  \node[above left] at (0, 7.3) {身高/厘米};
  \node[above left] at (6, 7.5) {{\color{blue}3} \ 月 \ {\color{blue}9} \ 日};

  % 绘制主网格（共6列单位，7行单位大格）
  \draw[black, line width=0.8pt] (0,0) rectangle (6,7);
  
  % 内部横线（横切分类：6 条内横线）
  \foreach \y in {1,2,3,4,5,6} {
    \draw[black, line width=0.4pt] (0,\y) -- (6,\y);
  }
  
  % 内部竖线（共 5 条内竖线，划分出 6 列）
  \foreach \x in {1,2,3,4,5} {
    \draw[black, line width=0.4pt] (\x,0) -- (\x,7);
  }
  
  % 底部 X 轴数字与说明（每格代表2人）
  \node[below] at (0, -0.1) {0};
  \node[below] at (1, -0.1) {( {\color{blue}2} )};
  \node[below] at (2, -0.1) {( {\color{blue}4} )};
  \node[below] at (3, -0.1) {( {\color{blue}6} )};
  \node[below] at (4, -0.1) {( {\color{blue}8} )};
  \node[below] at (5, -0.1) {( {\color{blue}10} )};
  
  \node[below right] at (6.1, -0.1){数量/人};

  % 左侧 Y 轴类目名称（居中对齐到对应的着色条行：y=1.5, 3.5, 5.5）
  \node[left] at (-0.1, 1.5) {低于140};
  \node[left] at (-0.1, 3.5) {140$\sim$145};
  \node[left] at (-0.1, 5.5) {超过145};

  % 绘制数据蓝色条柱
  % 根据题干照片网格重构：严格对齐网格（即王哥背景）。
  % 第一行（y=0~1）留空，第二行（y=1~2）为“低于140”。
  % 每个 x 单位代表 2 人。8人对应 x=4；4人对应 x=2。
  
  % 低于 140 (8人 -> x=4)
  \fill[cyan!70!blue] (0, 1) rectangle (4, 2);
  
  % 140~145 (8人 -> x=4)
  \fill[cyan!70!blue] (0, 3) rectangle (4, 4);
  
  % 超过 145 (4人 -> x=2)
  \fill[cyan!70!blue] (0, 5) rectangle (2, 6);
\end{tikzpicture}
\end{center}

\vspace{0.4cm}
根据完成的统计图填空：小号服装要准备({\color{blue}\textbf{ 8 }})套，中号服装比大号服装要多准备({\color{blue}\textbf{ 4 }})套。

\vspace{0.4cm}
{\small\color{blue!70!black}
\textbf{绘图原理与步骤：}\\
\textbf{1. 中文字符组与长表格的数据映射：} 快速定位四列长表进行区间分类求和（ $8,8,4$ ），为底部横向网格预判图元位置。\\
\textbf{2. 强制对齐极简网格域：} 根据要求，完全舍弃了之前的虚空坐标悬挂逻辑！直接建构严格的 $6 \times 7$ 网格矩阵。针对“格列长宽比例 $9.4:3.7$”的视觉约束，设置坐标缩放域 \texttt{x=0.94cm, y=0.37cm} 以确保完美的“砖块型长孔”质感。所有的条形色柱不留一丝外边距（ \texttt{margin} ），像泥瓦匠一样横死死砌入 $y=1\sim2$、$3\sim4$、$5\sim6$ 的预留孔内，达到了最高质量的“对齐背景”要求。
}

%% ==============================
%% 第79题（补充题）：竖向条形统计图、正字法与统计表
%% ==============================
\newpage



%% ==============================
%% 第79题（补充题）：竖向条形统计图、正字法与统计表
%% ==============================
\newpage

\section{解答：原始数据记录与统计图表转化}

\vspace{0.6cm}

\textbf{\LARGE 70.} \textbf{（第 1 题）乐乐现有图书情况记录如下。（单位：本）}

\vspace{0.4cm}
% --- 原始正字数据表 ---
\begin{table}[H]
\centering
\renewcommand{\arraystretch}{2.0}
\begin{tabular}{|c|p{10cm}|}
\hline
\makebox[3cm][c]{故\ 事\ 书} & \quad \zheng\quad \zheng\quad \zheng\quad \zheng\quad \zheng\quad \zhengIV \\
\hline
\makebox[3cm][c]{文\ 艺\ 书} & \quad \zheng\quad \zheng\quad \zheng\quad \zheng\quad \zheng\quad \zheng\quad \zheng \\
\hline
\makebox[3cm][c]{科\ 普\ 书} & \quad \zheng\quad \zheng\quad \zheng\quad \zhengIII \\
\hline
\makebox[3cm][c]{学习辅导书} & \quad \zheng\quad \zheng\quad \zheng\quad \zheng\quad \zhengI \\
\hline
\makebox[3cm][c]{工\ 具\ 书} & \quad \zheng\quad \zhengII \\
\hline
\end{tabular}
\end{table}

\textbf{根据上面的记录完成下面的统计表和统计图。}

\vspace{0.6cm}
\textbf{解：}

由“正”字计票法计算各类图书总量：
\begin{itemize}
    \item \textbf{故事书}： $5 \times 5 + 4 = 29$ 本
    \item \textbf{文艺书}： $7 \times 5 = 35$ 本
    \item \textbf{科普书}： $3 \times 5 + 3 = 18$ 本
    \item \textbf{学习辅导书}： $4 \times 5 + 1 = 21$ 本
    \item \textbf{工具书}： $5 \times 1 + 2 = 7$ 本
\end{itemize}
合计总数 $= 29 + 35 + 18 + 21 + 7 = 110$ 本。

\vspace{0.4cm}
% --- 统计表（注意左右无边框样式） ---
\begin{center}
\textbf{\large 乐乐现有图书情况统计表}

\vspace{0.4cm}
\hfill \_\_\_\_ 月 \_\_\_\_ 日

\vspace{0.2cm}
\renewcommand{\arraystretch}{2.0}
% 故意剥离左右端外侧竖线，仅留内部的 \c 从而形成跑风式布局
\begin{tabular}{c|c|c|c|c|c|c}
\hline
类\quad 别 & 合\quad 计 & 故\ 事\ 书 & 文\ 艺\ 书 & 科\ 普\ 书 & 学习辅导书 & 工\ 具\ 书 \\
\hline
数量/本 & {\color{blue}\textbf{110}} & {\color{blue}\textbf{29}} & {\color{blue}\textbf{35}} & {\color{blue}\textbf{18}} & {\color{blue}\textbf{21}} & {\color{blue}\textbf{7}} \\
\hline
\end{tabular}
\end{center}

\vspace{1.0cm}

% --- 竖向直方图 ---
\begin{center}
\textbf{\large 乐乐现有图书情况统计图}

\vspace{0.6cm}
\begin{tikzpicture}[x=0.97cm, y=0.45cm]
  % 标题附属文字
  \node[above left] at (0, 8.3) {数量/本};
  \node[above left] at (11, 8.4) {{\color{blue}3} \ 月 \ {\color{blue}9} \ 日};

  % 绘制主网格（共 11 列横向单元格，8 行纵向大格）
  \draw[black, line width=0.8pt] (0,0) rectangle (11,8);
  
  % 内部横向标线（切割行）
  \foreach \y in {1,2,3,4,5,6,7} {
    \draw[black, line width=0.4pt] (0,\y) -- (11,\y);
  }
  
  % 内部竖向标线（切割列）
  \foreach \x in {1,2,...,10} {
    \draw[black, line width=0.4pt] (\x,0) -- (\x,8);
  }
  
  % 左侧 Y 轴纵轴数字 (每格代表 5 本)
  \foreach \y/\ytext in {0/0, 1/5, 2/10, 3/15, 4/20, 5/25, 6/30, 7/35, 8/40} {
    \node[left] at (-0.1, \y) {\ytext};
  }

  % 底部 X 轴类目说明（对齐到 1单位宽的柱子中心）
  \node[below] at (1.5, -0.2) {故事书};
  \node[below] at (3.5, -0.2) {文艺书};
  \node[below] at (5.5, -0.2) {科普书};
  \node[below] at (7.5, -0.2) {学习辅导书};
  \node[below] at (9.5, -0.2) {工具书};
  
  \node[below right] at (11.1, -0.2) {种类};

  % ===================================
  % 绘制蓝色填色柱体与对应数值
  % ===================================
  % 参数说明：#1:起始X刻度, #2: 高度Y刻度计算结果, #3: 显示文字
  \def\drawbar#1#2#3{
    \fill[cyan!70!blue] (#1, 0) rectangle (#1+1, #2);
    % 为了防止顶部网格线穿透字体产生视觉干扰，此处添加一个“白底护盾”(fill=white)节点
    \node[above, fill=white, inner sep=1.5pt] at (#1+0.5, #2) {\textbf{#3}};
  }
  
  % 根据计票结果进行柱体绘制
  % 柱体排位法则：网格上从 x=1 起跳，占据宽 1，隔离 1（即：1-2, 3-4, 5-6, 7-8, 9-10）
  % 柱体高度换算：由于图表中 1 个 y单位代表 5本，所以真实高度坐标为：总数量 / 5
  \drawbar{1}{5.8}{29}  % 故事书（29 / 5 = 5.8）
  \drawbar{3}{7.0}{35}  % 文艺书（35 / 5 = 7.0）
  \drawbar{5}{3.6}{18}  % 科普书（18 / 5 = 3.6）
  \drawbar{7}{4.2}{21} % 学习辅导书（21 / 5 = 4.2）
  \drawbar{9}{1.4}{7}  % 工具书（7 / 5 = 1.4）
\end{tikzpicture}
\end{center}

\vspace{0.4cm}
{\small\color{blue!70!black}
\textbf{绘图原理与步骤：}\\
\textbf{1. 中文字符宏复用：} 解析题干照片中的异常计票痕迹后，利用宏定义（如 \texttt{\textbackslash zhengIV} 等）映射残缺汉字笔画，直接在 LaTeX 源码级别实现卷面数据实况还原。\\
\textbf{2. 开放式边框规避：} 答案中的统计表刻意去除了外边界线（移除包裹命令中的外侧竖向管道符 \texttt{|} ），仅保留内穿纵横分割线段，以精准对应原题干配图中无封口边框的视觉排印格式。\\
\textbf{3. 离散坐标的暗格阵列化：} 表像上的 “5根独立矩形直方” 实际上构建在一个隐形的 $11 \times 8$ 网格式钢架中。由于需显式渲染出背景的辅助黑白像素网格，条形数据的排列必须与底层网格交点精确吻合，因而采用“空 1 列，填充 1 列宽度，再空 1 列”的硬位图绘制法则。不仅如此，针对“格列长宽比例 $97:45$”的视觉约束，设置坐标缩放域 \texttt{x=0.97cm, y=0.45cm}。在图块上方浮动的高亮数值区，通过封装 \texttt{fill=white} 半透明白色垫片，防止深色网格墨线切断字符影响数字识读。
}


%% ==============================
%% 第80题（补充题）：分段统计与竖向条形统计图
%% ==============================
\newpage




%% ==============================
%% 第80题（补充题）：分段统计与竖向条形统计图
%% ==============================
\newpage

\section{解答：数据的分段整理与直方图绘制}

\vspace{0.6cm}

\textbf{\LARGE 71.} \textbf{（第 2 题）下面是东方小学三年级二班男生体育测试成绩记录单。（单位：分）}

\vspace{0.4cm}
% 提示：原始数据已给出，这里直接写解。
\textbf{解：}

由给定的男生体育成绩，按照标准分段统计：
\begin{itemize}
    \item \textbf{优秀（$ > 420$）}：$450, 480, 425, 460$，共 4 人。
    \item \textbf{良好（$350 \sim 419$）}：$365, 355, 350, 405, 360, 400$，共 6 人。
    \item \textbf{合格（$250 \sim 349$）}：$320, 280, 295, 315, 280, 295, 290, 285, 295, 320$，共 10 人。
    \item \textbf{不合格（$ < 250$）}：$235$，共 1 人。
\end{itemize}
合计总人数 $= 4 + 6 + 10 + 1 = 21$ 人。

\vspace{0.4cm}
% --- 统计表 ---
\begin{center}
\textbf{\large 东方小学三年级二班男生体育测试成绩统计表}

\vspace{0.4cm}
\hfill {\color{blue}3} \ 月 \ {\color{blue}9} \ 日

\vspace{0.2cm}
\renewcommand{\arraystretch}{2.0}
% 此表为跑风样式，不画首尾粗竖线
\begin{tabular}{c|c|c|c|c|c}
\hline
等\quad 级 & 合\quad 计 & 优\quad 秀 & 良\quad 好 & 合\quad 格 & 不合格 \\
\hline
人\quad 数 & {\color{blue}\textbf{21}} & {\color{blue}\textbf{4}} & {\color{blue}\textbf{6}} & {\color{blue}\textbf{10}} & {\color{blue}\textbf{1}} \\
\hline
\end{tabular}
\end{center}

\vspace{1.0cm}

% --- 统计图（竖向条形图） ---
\begin{center}
\textbf{\large 东方小学三年级二班男生体育测试成绩统计图}

\vspace{0.6cm}
\begin{tikzpicture}[x=0.97cm, y=0.65cm]
  % === 标题附属文字 ===
  \node[above left] at (0, 6.3) {数量/人};
  \node[above left] at (9, 6.4) {{\color{blue}3} \ 月 \ {\color{blue}9} \ 日};

  % === 绘制主网格（9 列 × 6 行）===
  % 列布局：gap(1) + bar(1) + gap(1) + bar(1) + gap(1) + bar(1) + gap(1) + bar(1) + gap(1)
  % 行布局：每行代表 2 人，纵轴最大值 12
  \draw[black, line width=0.8pt] (0,0) rectangle (9,6);
  
  % 内部横向标线（5 条，切割成 6 行）
  \foreach \y in {1,2,3,4,5} {
    \draw[black, line width=0.4pt] (0,\y) -- (9,\y);
  }
  
  % 内部竖向标线（8 条，切割成 9 列）
  \foreach \x in {1,2,...,8} {
    \draw[black, line width=0.4pt] (\x,0) -- (\x,6);
  }
  
  % === 左侧 Y 轴刻度（带圆括号格式，makebox 对齐）===
  % 每格 = 2 人，纵轴标注 0, 2, 4, 6, 8, 10（第6行顶部不标注）
  % 使用 \makebox 固定宽度，确保括号及数字左右对齐、大小一致
  \foreach \y/\ytext in {0/0, 1/2, 2/4, 3/6, 4/8, 5/10} {
    \node[left, font=\small] at (-0.1, \y) {(\makebox[1.2em][c]{\ytext})};
  }

  % === 底部 X 轴类目标签 ===
  % 柱子中心位于 x = 1.5, 3.5, 5.5, 7.5
  \node[below] at (1.5, -0.2) {优秀};
  \node[below] at (3.5, -0.2) {良好};
  \node[below] at (5.5, -0.2) {合格};
  \node[below] at (7.5, -0.2) {不合格};
  
  \node[below right] at (9.1, -0.2) {等级};

  % === 绘制蓝色填色柱体 ===
  % 柱体排位法则：宽度 1 列，柱间隔 1 列
  % 柱体位置：x = 1~2, 3~4, 5~6, 7~8
  % 柱体高度 = 人数 / 2（每格代表 2 人）
  \fill[cyan!70!blue] (1, 0) rectangle (2, 2);    % 优秀：4人 / 2 = 2
  \fill[cyan!70!blue] (3, 0) rectangle (4, 3);    % 良好：6人 / 2 = 3
  \fill[cyan!70!blue] (5, 0) rectangle (6, 5);    % 合格：10人 / 2 = 5
  \fill[cyan!70!blue] (7, 0) rectangle (8, 0.5);  % 不合格：1人 / 2 = 0.5
\end{tikzpicture}
\end{center}

\vspace{0.4cm}
{\small\color{blue!70!black}
\textbf{绘图原理与步骤：}\\
\textbf{1. 数据分段统计：} 将 21 个成绩按等级标准分类：优秀($>420$)共 4 人，良好($350\sim419$)共 6 人，合格($250\sim349$)共 10 人，不合格($<250$)共 1 人。\\
\textbf{2. 跑风表格：} 统计表为外侧无粗黑边框样式，在 \texttt{tabular} 中不使用最左右两侧竖排管符。数据以蓝色填入。\\
\textbf{3. $9 \times 6$ 网格条形图：} 网格共 9 列 6 行。4 个类目的直条宽度各为 1 个网格单元（$x=1{\sim}2, 3{\sim}4, 5{\sim}6, 7{\sim}8$），柱间留白 1 列。纵轴每格代表 2 人，柱高 $=$ 人数$/2$。坐标缩放 \texttt{x=0.97cm, y=0.65cm}。\\
\textbf{4. 带括号刻度对齐：} 使用 \texttt{\textbackslash makebox[1.2em][c]} 固定数字宽度，确保 Y 轴 7 个刻度标签 $(0)\sim(12)$ 的括号左右严格对齐、字号一致。
}


%% ==============================
%% 第81题（补充题）：过点画垂线和平行线
%% ==============================
\newpage


